\documentclass[11pt,a4paper]{article}
\usepackage[utf8]{inputenc}
\usepackage[T1]{fontenc}
\usepackage[brazil]{babel}
\usepackage{lmodern}
\usepackage{microtype}
\usepackage{amsmath,amssymb,amsthm,mathtools}
\usepackage[margin=2.2cm]{geometry}
\usepackage{enumitem}
\usepackage{booktabs,tabularx,multirow,array}
\usepackage[dvipsnames,table]{xcolor}
\usepackage[most]{tcolorbox}
\usepackage{tikz}
\usetikzlibrary{arrows.meta,positioning,shapes.geometric,calc,decorations.pathreplacing}
\usepackage{forest}
\usepackage{graphicx}
\usepackage[noend]{algpseudocode}
\usepackage[colorlinks=true,linkcolor=MidnightBlue,urlcolor=MidnightBlue]{hyperref}

% ------------------------------------------------------------------
% Macros
% ------------------------------------------------------------------
\newcommand{\N}{\mathbb{N}}
\newcommand{\R}{\mathbb{R}}
\newcommand{\bal}{\operatorname{bal}}
\newcommand{\rank}{\operatorname{rank}}
\newcommand{\Find}{\textsc{Find}}
\newcommand{\Union}{\textsc{Union}}
\newcommand{\MakeSet}{\textsc{Make-Set}}
\newcommand{\lgs}{\log^{*}}
\DeclarePairedDelimiter{\floor}{\lfloor}{\rfloor}
\DeclarePairedDelimiter{\ceil}{\lceil}{\rceil}
\DeclarePairedDelimiter{\abs}{\lvert}{\rvert}
\newcommand{\am}{$^{*}$}   % marca de amortizado

% ------------------------------------------------------------------
% Pseudocódigo em português (algpseudocode)
% ------------------------------------------------------------------
\algrenewcommand\algorithmicif{\textbf{se}}
\algrenewcommand\algorithmicthen{\textbf{então}}
\algrenewcommand\algorithmicelse{\textbf{senão}}
\algrenewcommand\algorithmicwhile{\textbf{enquanto}}
\algrenewcommand\algorithmicdo{\textbf{faça}}
\algrenewcommand\algorithmicfor{\textbf{para}}
\algrenewcommand\algorithmicreturn{\textbf{retorna}}
\algrenewcommand\algorithmicprocedure{\textbf{Procedimento}}
\algrenewcommand\algorithmicfunction{\textbf{Função}}
\algrenewcommand\algorithmicrepeat{\textbf{repita}}
\algrenewcommand\algorithmicuntil{\textbf{até que}}
\algrenewcommand{\algorithmiccomment}[1]{\hfill{\color{gray}\small\% #1}}
\algnewcommand\algorithmicelsif{\textbf{senão se}}
\algdef{C}[IF]{IF}{ElsIf}[1]{\algorithmicelsif\ #1\ \algorithmicthen}
\newcommand{\ptr}{\mathbin{\uparrow}}   % pt ↑ .campo
\newcommand{\nil}{\lambda}

% ------------------------------------------------------------------
% Ambientes
% ------------------------------------------------------------------
\theoremstyle{plain}
\newtheorem{teo}{Teorema}[section]
\newtheorem{prop}[teo]{Proposição}
\newtheorem{lema}[teo]{Lema}
\newtheorem{cor}[teo]{Corolário}
\theoremstyle{definition}
\newtheorem{defi}[teo]{Definição}
\newtheorem{exemplo}[teo]{Exemplo}
\theoremstyle{remark}
\newtheorem{obs}[teo]{Observação}

% Pseudocódigo em caixa (sem float)
\newtcolorbox{pseudo}[1]{enhanced, breakable, colback=gray!3, colframe=gray!45, colbacktitle=gray!22, coltitle=black,
  fonttitle=\bfseries\small, title={#1}, left=3pt,right=3pt,top=2pt,bottom=2pt, fontupper=\small}
% Pontos-chave
\newtcolorbox{chave}[1][Pontos-chave]{enhanced, breakable, colback=MidnightBlue!4,
  colframe=MidnightBlue!70!black, fonttitle=\bfseries, title={#1}, left=4pt,right=4pt,top=3pt,bottom=3pt}
% Como cai na prova
\newtcolorbox{prova}[1][Como isso cai na AV1]{enhanced, breakable, colback=ForestGreen!5,
  colframe=ForestGreen!70!black, fonttitle=\bfseries, title={#1}, left=4pt,right=4pt,top=3pt,bottom=3pt}
% Armadilhas
\newtcolorbox{cuidado}[1][Armadilhas comuns]{enhanced, breakable, colback=BrickRed!4,
  colframe=BrickRed!75!black, fonttitle=\bfseries, title={#1}, left=4pt,right=4pt,top=3pt,bottom=3pt}
% Complemento (fora das aulas)
\newtcolorbox{complemento}[1][]{enhanced, breakable, colback=orange!4, colframe=orange!70!black,
  fonttitle=\bfseries, title={Complemento#1}, left=4pt,right=4pt,top=3pt,bottom=3pt}

% Árvores (forest): nó circular com balanço opcional
\forestset{
  bst/.style={for tree={circle, draw, minimum size=1.5em, inner sep=1pt, font=\small,
      s sep=6mm, l sep=4mm, edge={thick}}},
  bb/.style={label={[font=\scriptsize,text=BrickRed,inner sep=1pt]right:{$#1$}}},
  vazio/.style={phantom},
  tri/.style={draw, isosceles triangle, shape border rotate=90, isosceles triangle apex angle=60,
      inner sep=0.5pt, minimum height=7mm, font=\scriptsize, anchor=apex, child anchor=apex},
  novo/.style={fill=yellow!40},
  desreg/.style={fill=BrickRed!25},
}

\setlist{itemsep=1.5pt,topsep=3pt}
\numberwithin{equation}{section}
\renewcommand{\arraystretch}{1.15}

\title{\textbf{MC23 -- Algoritmos}\\[4pt]
\large Resumo detalhado para a AV1\\
(IMPA Tech, 2026.2 -- Prof.\ Uéverton Souza)}
\author{Baseado nos slides, aulas manuscritas, monitorias, listas, Testes 1 e 2;\\
complementado com Szwarcfiter \& Markenzon (\emph{Estruturas de Dados e seus Algoritmos})\\
e Cormen et al.\ (\emph{Introduction to Algorithms}, 3.\ ed.)}
\date{AV1: terça-feira, 22 de setembro de 2026}

\begin{document}
\maketitle

\begin{chave}[Como usar este resumo]
\begin{itemize}
  \item Ele foi escrito para quem \textbf{já estudou a matéria} e quer se localizar rápido: cada seção começa com o essencial, depois vêm as demonstrações e os detalhes que costumam ser cobrados.
  \item \textcolor{ForestGreen!70!black}{\textbf{Caixas verdes}}: como o tema cai na AV1 (o formato da prova foi anunciado em aula --- ver abaixo).
  \item \textcolor{BrickRed!75!black}{\textbf{Caixas vermelhas}}: armadilhas e erros frequentes (inclusive erros encontrados em slides/monitorias).
  \item \textcolor{MidnightBlue}{\textbf{Caixas azuis}}: resumo do que precisa estar na ponta da língua.
  \item \textcolor{orange!70!black}{\textbf{Caixas laranjas}}: complementos (livros/pesquisas) que não foram o foco das aulas.
  \item Convenção de heap do curso (Szwarcfiter): \textbf{heap de máximo, vetor indexado a partir de 1}, pai de $i$ é $\floor{i/2}$, filhos $2i$ e $2i+1$. Balanço AVL: $\bal(v)=h_D(v)-h_E(v)$.
  \item Todos os exemplos numéricos (AVL, heaps, radix, subvetor máximo, inversões) foram \emph{conferidos por simulação}.
\end{itemize}
\end{chave}

\begin{prova}[Formato da AV1 (anunciado na aula de revisão de 18/09)]
\begin{center}
\begin{tabularx}{\linewidth}{@{}lXl@{}}
\toprule
\textbf{Questão} & \textbf{O que será pedido} & \textbf{Onde estudar} \\
\midrule
Q1 & \textbf{Árvore AVL}: saber como a estrutura funciona, resolver uma tarefa nela, manipular (inserir/remover, rotações, balanços). & Seção \ref{sec:avl} \\
Q2 & \textbf{Union-Find} (a mais fácil): diferentes formas de implementar e em que situações cada uma é melhor. & Seção \ref{sec:uf} (Tabela \ref{tab:uf}) \\
Q3 & \textbf{Heap}: saber \emph{demonstrar} propriedades de heap. & §\ref{sec:heapprops} (provas); depois §\ref{sec:heap}, §\ref{sec:heapav} \\
Q4 & Dado um código genérico: qual \textbf{estrutura de dados} usar e qual \textbf{ordenação} usar. & Seções \ref{sec:escolhaheap}, \ref{sec:escolhasort}, \ref{sec:escolhageral} \\
Q5 & \textbf{Divisão e conquista}: dado um problema, escrever o \textbf{pseudocódigo}. & Seção \ref{sec:dc} \\
\bottomrule
\end{tabularx}
\end{center}
\end{prova}

\tableofcontents
\newpage

% ================================================================
\section*{Mapa: aulas e materiais $\to$ seções}
\addcontentsline{toc}{section}{Mapa: aulas e materiais $\to$ seções}

\begin{center}
\small
\begin{tabularx}{\linewidth}{@{}llX l@{}}
\toprule
\textbf{Aula} & \textbf{Data} & \textbf{Conteúdo} & \textbf{Seção}\\
\midrule
1 & 11/08 & Complexidade, notações assintóticas, limites superior/inferior, polinomial $\times$ exponencial & \ref{sec:complex}\\
2 & 12/08 & Listas de prioridades e heaps (binário), heapsort & \ref{sec:heap}\\
3 & 18/08 & Heaps binários (resumo de custos) e heaps binomiais & \ref{sec:heap}, \ref{sec:heapav}\\
4 & 21/08 & Heaps de Fibonacci e complexidade amortizada & \ref{sec:amort}, \ref{sec:heapav}\\
5 & 25/08 & Union-Find: listas, florestas, rank, compressão, análise $O(m\lgs n)$ & \ref{sec:uf}\\
6 & 28/08 & Árvores binárias de busca e AVL & \ref{sec:avl}\\
7 & 01/09 & Ordenação por comparação (heap/merge/quick), limite inferior $\Omega(n\log n)$ & \ref{sec:sortcomp}\\
8 & 04/09 & Ordenação linear: counting, radix, bucket (e American Flag) & \ref{sec:sortlin}\\
9 & 08/09 & Seleção: algoritmo BFPRT (mediana das medianas) & \ref{sec:selecao}\\
10 & 11/09 & Divisão e conquista; contagem de inversões & \ref{sec:dc}\\
11 & 15/09 & \textbf{Teste 2} (ordenação e D\&C) & \ref{sec:testes}\\
12--13 & 16 e 18/09 & Exercícios; revisão (AVL, D\&C); formato da AV1 & ---\\
\midrule
Monitorias & & Heaps binomiais/Fibonacci (26/08), Union-Find, AVL (aula prática 4), Revisão Heap/UF/AVL & \ref{sec:heapav}, \ref{sec:uf}, \ref{sec:avl}, \ref{sec:banco}\\
Pesquisas & até 19/09 & Árvores rubro-negras; Tabelas hash & \ref{sec:rn}, \ref{sec:hash}\\
Listas & & Heaps (binário/binomial/Fibonacci); Union-Find e AVL & \ref{sec:banco}\\
\bottomrule
\end{tabularx}
\end{center}

\noindent\emph{Obs.:} os slides de Programação Dinâmica, Gulosos, Matroides, AGM e Grafos pertencem à parte da AV2 no andamento real do curso; só PD aparece aqui, como complemento curto (§\ref{sec:pd}).

\newpage
% ================================================================
\section*{Cola: complexidades de tudo}
\addcontentsline{toc}{section}{Cola: complexidades de tudo}

\begin{center}
\small
\textbf{Listas de prioridades} ($n$ elementos; \am\,= amortizado; heap de máximo)\\[2pt]
\resizebox{\linewidth}{!}{\begin{tabular}{@{}lccccc@{}}
\toprule
Operação & Lista não ord. & Lista ordenada & Heap binário & Heap binomial & Heap de Fibonacci\\
\midrule
Seleção (máx.) & $\Theta(n)$ & $\Theta(1)$ & $O(1)$ & $O(\log n)$ ou $O(1)$ c/ ponteiro & $O(1)$\\
Inserção & $\Theta(1)$ & $\Theta(n)$ & $O(\log n)$ & $O(\log n)$; $O(1)$ esperado/\am & $O(1)$\am\\
Remoção do máx. & $\Theta(n)$ & $\Theta(1)$ & $O(\log n)$ & $O(\log n)$ & $O(\log n)$\am\\
Aumentar prioridade & $\Theta(n)$ & $\Theta(n)$ & $O(\log n)$ & $O(\log n)$ & $O(1)$\am\\
Diminuir prioridade & $\Theta(n)$ & $\Theta(n)$ & $O(\log n)$ & $O(\log n)$ & $O(\log n)$\am\\
Construção & $\Theta(n)$ & $\Theta(n\log n)$ & $O(n)$ & $O(n)$ & $O(n)$\\
União (fusão) & $O(1)$ encad. & $O(n)$ & $O(n)$ & $O(\log n)$ & $O(1)$\am\\
\bottomrule
\end{tabular}}

\medskip
\textbf{Union-Find} ($n$ elementos, $m$ operações)\\[2pt]
\resizebox{\linewidth}{!}{\begin{tabular}{@{}lccc@{}}
\toprule
Implementação & \MakeSet & \Find & \Union\\
\midrule
Listas encadeadas c/ ponteiro ao representante (união pelo tamanho) & $O(1)$ & $O(1)$ & $O(\log n)$\am\\
Floresta ingênua & $O(1)$ & $O(n)$ & $O(1)$ (dadas as raízes)\\
Floresta + união por tamanho/rank & $O(1)$ & $O(\log n)$ & $O(1)$ (dadas as raízes)\\
Floresta + rank + compressão de caminhos & $O(1)$ & \multicolumn{2}{c}{$O(m\lgs n)$ no total (justo: $O(m\,\alpha(m,n))$)}\\
\bottomrule
\end{tabular}}

\medskip
\textbf{Árvores de busca}\quad ABB: busca/inserção/remoção $\Theta(h)$, com $h$ entre $\floor{\log n}+1$ e $n$.\quad AVL: tudo $O(\log n)$ no pior caso; inserção faz $\le 1$ rotação (simples ou dupla); remoção pode fazer $O(\log n)$ rotações.

\medskip
\textbf{Ordenação}\\[2pt]
\resizebox{\linewidth}{!}{\begin{tabular}{@{}lcccccc@{}}
\toprule
Algoritmo & Pior caso & Caso médio & Melhor caso & Memória extra & Estável? & In-place?\\
\midrule
Inserção / Bolha & $\Theta(n^2)$ & $\Theta(n^2)$ & $\Theta(n)$ & $O(1)$ & sim & sim\\
Mergesort & $\Theta(n\log n)$ & $\Theta(n\log n)$ & $\Theta(n\log n)$ & $\Theta(n)$ & sim & não\\
Heapsort & $\Theta(n\log n)$ & $\Theta(n\log n)$ & $\Theta(n\log n)$ & $O(1)$ & não & sim\\
Quicksort & $\Theta(n^2)$ & $\Theta(n\log n)$ & $\Theta(n\log n)$ & $O(\log n)$ pilha (média) & não & sim\\
Counting sort & $\Theta(n+k)$ & $\Theta(n+k)$ & $\Theta(n+k)$ & $\Theta(n+k)$ & sim & não\\
Radix sort (LSD) & $\Theta(d(n+b))$ & idem & idem & $\Theta(n+b)$ & sim & não\\
Bucket sort & $\Theta(n\log n)$ ou $\Theta(n^2)$ & $\Theta(n)$ (unif.) & $\Theta(n)$ & $\Theta(n)$ & depende & não\\
\bottomrule
\end{tabular}}

\medskip
\textbf{Outros}\quad Limite inferior p/ ordenação por comparação: $\Omega(n\log n)$.\quad Seleção ($k$-ésimo menor): BFPRT $O(n)$ pior caso; quickselect $O(n)$ esperado, $O(n^2)$ pior.\quad Contagem de inversões e subvetor de soma máxima (D\&C): $O(n\log n)$.\quad Strassen: $O(n^{\log_2 7})\approx O(n^{2{,}807})$.
\end{center}

\newpage
% ================================================================
\section{Complexidade de algoritmos}\label{sec:complex}
% ================================================================

\subsection{Medidas de complexidade}
Seja $A$ um algoritmo e $\{E_1,\dots,E_m\}$ todas as entradas possíveis de tamanho $n$; $T_i$ é o número de passos de $A$ sobre $E_i$.
\begin{itemize}
  \item \textbf{Pior caso:} $\max_i T_i$ \quad(é a medida padrão do curso).
  \item \textbf{Caso médio:} $\sum_i p_i T_i$, onde $p_i$ é a probabilidade de $E_i$.
  \item \textbf{Melhor caso:} $\min_i T_i$.
  \item \textbf{Amortizada:} custo de uma \emph{sequência} de operações dividido pelo número de operações, no pior caso da sequência (§\ref{sec:amort}). Não é média probabilística!
\end{itemize}

\paragraph{Tamanho da entrada.} É o comprimento da codificação da entrada (em bits/símbolos). Isso importa com números: um algoritmo é \textbf{pseudopolinomial} se é polinomial quando os números da entrada são codificados em \emph{unário}. Exemplo do curso: o \emph{counting sort} é $O(n+k)$; se $k=\Theta(2^n)$ ele é exponencial no tamanho da entrada (um número $k$ ocupa só $\log k$ bits).

\subsection{Notações assintóticas}
Sejam $f,g$ funções (assintoticamente) não negativas.
\begin{center}
\begin{tabular}{@{}llll@{}}
\toprule
Notação & Definição & Via limite (se existir) & Leitura\\
\midrule
$f=O(g)$ & $\exists c>0,n_0:\ f(n)\le c\,g(n)\ \forall n\ge n_0$ & $\lim f/g<\infty$ & ``$\le$''\\
$f=\Omega(g)$ & $\exists c>0,n_0:\ f(n)\ge c\,g(n)\ \forall n\ge n_0$ & $\lim f/g>0$ & ``$\ge$''\\
$f=\Theta(g)$ & $f=O(g)$ e $f=\Omega(g)$ & $\lim f/g=c\in(0,\infty)$ & ``$=$''\\
$f=o(g)$ & $\forall c>0\ \exists n_0:\ f(n)<c\,g(n)$ & $\lim f/g=0$ & ``$<$'' (não justo)\\
$f=\omega(g)$ & $\forall c>0\ \exists n_0:\ f(n)>c\,g(n)$ & $\lim f/g=\infty$ & ``$>$'' (não justo)\\
\bottomrule
\end{tabular}
\end{center}
Exemplo dos slides: $2n^2+1$ é $O(n^3)$, $O(n^2)$, não é $O(n)$; é $\Omega(n)$, $\Omega(n^2)$, não é $\Omega(n^3)$; é $\Theta(n^2)$.
$O(g)$ também denota o \emph{conjunto} das funções $f$ com $f=O(g)$ (escreve-se $f\in O(g)$).

\paragraph{Propriedades úteis.} Transitividade; $f=O(g)\iff g=\Omega(f)$; $O(f)+O(g)=O(\max\{f,g\})$; $\log_a n=\Theta(\log_b n)$ (a base do log não importa); $\log(n!)=\Theta(n\log n)$; $(\log n)^c=o(n^\varepsilon)$ e $n^c=o(a^n)$ para $a>1$.
\[
1\ <\ \log\log n\ <\ \log n\ <\ \sqrt n\ <\ n\ <\ n\log n\ <\ n^2\ <\ n^3\ <\ n^{\log n}\ <\ 2^n\ <\ n!\ <\ n^n .
\]

\subsection{Complexidade de problemas $\times$ de algoritmos}
\begin{defi}\leavevmode
\begin{itemize}
  \item Um \textbf{limite superior} para o problema $\Pi$ é uma função $f(n)$ tal que \emph{existe} algum algoritmo que resolve $\Pi$ com no máximo $f(n)$ passos no pior caso. (A complexidade de pior caso de qualquer algoritmo para $\Pi$ é um limite superior para $\Pi$.)
  \item Um \textbf{limite inferior} para $\Pi$ é uma função $f(n)$ tal que \emph{todo} algoritmo para $\Pi$ precisa de pelo menos $f(n)$ passos no pior caso.
  \item Um algoritmo é \textbf{ótimo} se sua complexidade de pior caso é igual a um limite inferior do problema (ex.: mergesort e heapsort para ordenação por comparação).
\end{itemize}
\end{defi}
As notações são só sobre funções: nada impede dizer ``o pior caso de $A$ é $\Omega(n^3)$'' ou ``o limite inferior de $\Pi$ é $O(\log n)$''.

\paragraph{Polinomial $\times$ exponencial.} Consenso: um algoritmo é útil na prática se é \textbf{polinomial} ($O(n^c)$; classe $P$ = problemas com algoritmo polinomial). Exponenciais: $2^{O(n)}$, $k^n$, $n!$, $2^{n^2}$, $n^{\log n}$. Polinomiais aproveitam melhor o avanço tecnológico: com máquina $10\times$ mais rápida, o maior $n$ resolvido em um dia é \emph{multiplicado} por uma constante ($n^2$: $\times 3{,}16$), enquanto para $2^n$ só \emph{soma} uma constante ($40\to 43$).

% ------------------------------------------------------------------
\subsection{Recorrências}\label{sec:recorrencias}
Algoritmos recursivos (D\&C) têm tempo dado por recorrências $T(n)=a\,T(n/b)+f(n)$: $a$ subproblemas de tamanho $n/b$ e custo $f(n)$ para dividir e combinar.

\paragraph{Árvore de recursão.} Some o custo por nível. Ex.: $T(n)=2T(n/2)+cn$: há $\log_2 n$ níveis, cada um custando $cn$ $\Rightarrow T(n)=\Theta(n\log n)$ (mergesort, contagem de inversões, subvetor máximo).

\paragraph{Substituição (indução).} Chuta-se $T(n)\le c\,g(n)$ e prova-se por indução forte (é assim que o slide prova $T(n)\le 20cn$ na BFPRT, §\ref{sec:selecao}).

\begin{teo}[Teorema Mestre -- Cormen §4.5]
Seja $T(n)=a\,T(n/b)+f(n)$ com $a\ge1$, $b>1$, e seja $E=\log_b a$ (o ``expoente crítico'').
\begin{enumerate}
  \item Se $f(n)=O(n^{E-\varepsilon})$ para algum $\varepsilon>0$: $T(n)=\Theta(n^{E})$ \quad(as folhas dominam).
  \item Se $f(n)=\Theta(n^{E})$: $T(n)=\Theta(n^{E}\log n)$ \quad(todos os níveis pesam igual).
  \item Se $f(n)=\Omega(n^{E+\varepsilon})$ e $a f(n/b)\le c f(n)$ para algum $c<1$: $T(n)=\Theta(f(n))$ \quad(a raiz domina).
\end{enumerate}
\end{teo}
\begin{center}
\small
\begin{tabular}{@{}lll@{}}
\toprule
Recorrência & Exemplo & Solução\\
\midrule
$T(n)=T(n/2)+O(1)$ & busca binária & $\Theta(\log n)$\\
$T(n)=2T(n/2)+O(1)$ & máximo por D\&C & $\Theta(n)$\\
$T(n)=2T(n/2)+\Theta(n)$ & mergesort, inversões, subvetor máximo & $\Theta(n\log n)$\\
$T(n)=8T(n/2)+\Theta(n^2)$ & produto de matrizes por blocos & $\Theta(n^3)$\\
$T(n)=7T(n/2)+\Theta(n^2)$ & Strassen & $\Theta(n^{\log_2 7})\approx\Theta(n^{2{,}807})$\\
$T(n)=T(n/5)+T(3n/4)+\Theta(n)$ & BFPRT & $\Theta(n)$ (pois $\tfrac15+\tfrac34<1$)\\
$T(n)=T(n-1)+\Theta(n)$ & quicksort no pior caso & $\Theta(n^2)$\\
$T(n)=T(9n/10)+T(n/10)+\Theta(n)$ & quicksort ``desbalanceado fixo'' & $\Theta(n\log n)$\\
\bottomrule
\end{tabular}
\end{center}

\begin{teo}[D\&C $c$-balanceado -- slide de D\&C]
Se um algoritmo divide a instância em $k$ partes de tamanho no máximo $n/c$ ($c>1$), com partes de tamanhos que diferem por constante, e gasta tempo linear para dividir e combinar, então roda em $O(n\log n)$.
\end{teo}
\begin{proof}[Ideia]
Na árvore de recursão, os subproblemas de um mesmo nível são disjuntos, então cada nível custa $O(n)$; como o tamanho cai por um fator $c$ a cada nível, há $O(\log_c n)$ níveis.
\end{proof}

\paragraph{Somatórios que aparecem.} $\sum_{i=1}^n i=\frac{n(n+1)}2$; $\sum_{k\ge0}x^k=\frac1{1-x}$ ($|x|<1$); $\sum_{k\ge0}\frac{k}{2^k}=2$; $\sum_{k=0}^{\log n}\frac{n}{2^k}\le 2n$; $\sum_{i=1}^n\frac1i=\Theta(\log n)$; $\log(n!)=\sum\log i=\Theta(n\log n)$.

% ------------------------------------------------------------------
\subsection{Complexidade amortizada}\label{sec:amort}
A análise amortizada considera uma \textbf{sequência} de operações: cada operação individual pode ser cara, mas o custo total da sequência (no pior caso) é pequeno. \textbf{Não envolve probabilidade.}

\paragraph{Visão do banqueiro (créditos).} Cada operação recebe um número fixo de créditos (seu custo amortizado). Operações baratas ``poupam'' créditos, que pagam operações caras depois. Se o saldo nunca fica negativo, o custo real total $\le$ soma dos custos amortizados.

\paragraph{Visão do físico (potencial).} Escolhe-se $\Phi(E)\ge0$ para cada configuração $E$. Custo amortizado: $a_i=t_i+\Phi(E_i)-\Phi(E_{i-1})$. Somando (telescópica):
\[
\sum_{i=1}^m t_i=\sum_{i=1}^m a_i-\Phi_m+\Phi_0\ \le\ \sum_{i=1}^m a_i\quad\text{se }\Phi_m\ge\Phi_0 .
\]

\begin{exemplo}[Pilha com multi-desempilhamento -- Szwarcfiter §8.2]
Cada operação faz $k\ge0$ \emph{retira} e depois um \emph{insere}; no pior caso uma operação custa $\Theta(n)$. Com $\Phi=$ número de elementos na pilha: $a_i=(1+k)+(p-k+1)-p=2$. Logo $m$ operações custam $\le 2m=O(m)$.
\end{exemplo}
\begin{exemplo}[Contador binário]
Incrementar um contador de $k$ bits pode trocar $k$ bits, mas $m$ incrementos a partir de $0$ trocam no total $\le 2m$ bits: com $\Phi=$ número de bits $1$, um incremento que zera $t$ bits e liga um tem $a=(t+1)+(1-t)=2$. É exatamente o argumento da \textbf{inserção em heap binomial} (§\ref{sec:binomial}): inserir $=$ somar 1 em binário.
\end{exemplo}

\begin{chave}
``Amortizado'' $\ne$ ``caso médio''. Amortizado é \emph{garantia de pior caso para a sequência}; uma operação isolada ainda pode ser cara. Isso é decisivo na escolha entre heap binomial e de Fibonacci quando o sistema exige \textbf{previsibilidade de cada operação individual} (Teste 1, Q1).
\end{chave}

% ================================================================
\section{Listas de prioridades e heaps binários}\label{sec:heap}
% ================================================================

\subsection{Lista de prioridades}
Tabela em que cada dado tem uma \textbf{prioridade}; operações: (1) \emph{seleção} do de maior prioridade; (2) \emph{inserção}; (3) \emph{remoção} do de maior prioridade; (4) \emph{alteração} de prioridade; (5) \emph{construção} a partir de um conjunto inicial. (Fusão entra com os heaps avançados.)

\begin{center}
\small
\begin{tabular}{@{}lccc@{}}
\toprule
Operação & Lista \textbf{não} ordenada & Lista ordenada & Heap\\
\midrule
seleção & $\Theta(n)$ & $\Theta(1)$ & $O(1)$\\
inserção & $\Theta(1)$ & $\Theta(n)$ & $O(\log n)$\\
remoção & $\Theta(n)$ & $\Theta(1)$ & $O(\log n)$\\
alteração & $\Theta(n)$ & $\Theta(n)$ & $O(\log n)$\\
construção & $\Theta(n)$ & $\Theta(n\log n)$ & $O(n)$\\
\bottomrule
\end{tabular}
\end{center}
As duas listas têm alguma operação de atualização $\Omega(n)$; o heap equilibra tudo em $O(\log n)$.

\subsection{Definição e representação}
\begin{defi}[Heap -- convenção do curso]
Um \textbf{heap} (de máximo) é uma lista linear $s_1,\dots,s_n$ tal que
\[
s_i\le s_{\floor{i/2}}\qquad\text{para todo }1<i\le n .
\]
(Heap de \emph{mínimo}: $s_i\ge s_{\floor{i/2}}$.)
\end{defi}
Visualiza-se como uma \textbf{árvore binária completa} numerada por níveis, da esquerda para a direita: o nó $i$ tem pai $\floor{i/2}$, filho esquerdo $2i$ e direito $2i+1$ (se $\le n$). A condição diz: \emph{cada nó tem prioridade $\ge$ à dos filhos}. Não há ponteiros: é só um vetor.

\begin{center}
\begin{forest} bst
[95,label={[font=\tiny]above:1}
  [60,label={[font=\tiny]left:2} [39,label={[font=\tiny]left:4} [33,label={[font=\tiny]left:8}] [,vazio]] [28,label={[font=\tiny]right:5}]]
  [78,label={[font=\tiny]right:3} [66,label={[font=\tiny]left:6}] [70,label={[font=\tiny]right:7}]]]
\end{forest}
\qquad
\raisebox{3em}{\begin{tabular}{c|c|c|c|c|c|c|c}
1&2&3&4&5&6&7&8\\\hline 95&60&78&39&28&66&70&33
\end{tabular}}
\end{center}

\paragraph{Fatos de índice} (usados em todas as provas):
\begin{itemize}
  \item o nó $i$ está no nível $\floor{\log_2 i}+1$ (raiz no nível 1); a árvore tem $\floor{\log_2 n}+1$ níveis;
  \item as \textbf{folhas} são exatamente as posições $\floor{n/2}+1,\dots,n$ (pois $i$ tem filho sse $2i\le n$); há $\ceil{n/2}$ folhas;
  \item convenção de altura do Szwarcfiter/slides: \textbf{folhas têm altura 1}.
\end{itemize}

\subsection{Subir e descer}
\begin{minipage}[t]{0.47\linewidth}
\begin{pseudo}{subir$(i)$ -- após \emph{aumentar} a prioridade de $i$}
\begin{algorithmic}
\Procedure{subir}{$i$}
  \State $j := \floor{i/2}$
  \If{$j\ge1$}
    \If{$T[i].chave > T[j].chave$}
      \State $T[i]\Leftrightarrow T[j]$
      \State \Call{subir}{$j$}
    \EndIf
  \EndIf
\EndProcedure
\end{algorithmic}
\end{pseudo}
\end{minipage}\hfill
\begin{minipage}[t]{0.5\linewidth}
\begin{pseudo}{descer$(i,n)$ -- após \emph{diminuir} a prioridade de $i$}
\begin{algorithmic}
\Procedure{descer}{$i,n$}
  \State $j := 2i$
  \If{$j\le n$}
    \If{$j<n$ \textbf{e} $T[j+1].chave > T[j].chave$}
      \State $j := j+1$ \Comment{maior filho}
    \EndIf
    \If{$T[i].chave < T[j].chave$}
      \State $T[i]\Leftrightarrow T[j]$
      \State \Call{descer}{$j,n$}
    \EndIf
  \EndIf
\EndProcedure
\end{algorithmic}
\end{pseudo}
\end{minipage}

\medskip
Ambos percorrem um caminho da árvore: $O(\log n)$. \textbf{Descer escolhe o maior filho} (se trocasse com o menor, o maior filho ficaria abaixo de um valor menor).

\subsection{Inserção, remoção, alteração}
\begin{pseudo}{Inserção, remoção do máximo e alteração de prioridade}
\begin{algorithmic}
\Procedure{inserir}{$novo$} \Comment{$M$ = capacidade}
  \If{$n<M$} \State $T[n+1]:=novo$;\ $n:=n+1$;\ \Call{subir}{$n$}
  \Else\ \textbf{overflow}
  \EndIf
\EndProcedure
\Procedure{remover}{}
  \If{$n\ne0$} \State agir$(T[1])$;\ $T[1]:=T[n]$;\ $n:=n-1$;\ \Call{descer}{$1,n$}
  \Else\ \textbf{underflow}
  \EndIf
\EndProcedure
\Procedure{alterar}{$i, p$}
  \State $antiga := T[i].chave$;\ $T[i].chave := p$
  \If{$p > antiga$} \State \Call{subir}{$i$} \Comment{aumentou $\Rightarrow$ sobe}
  \Else \State \Call{descer}{$i,n$} \Comment{diminuiu $\Rightarrow$ desce}
  \EndIf
\EndProcedure
\end{algorithmic}
\end{pseudo}

\begin{cuidado}[Erro clássico (monitoria de revisão, Q5)]
Um aluno escreveu: ``$H[i]:=p$; \textbf{se} $H[i]>H[\floor{i/2}]$ \textbf{então} Descer$(i)$ \textbf{senão} Subir$(i)$''. Está \textbf{invertido}: se o novo valor ficou maior que o pai, ele viola a propriedade \emph{com o pai} e precisa \textbf{subir}; se não ficou maior que o pai, a única violação possível é com os filhos, e ele precisa \textbf{descer}. Versão correta: \textbf{se} $H[i]>H[\floor{i/2}]$ (com $i>1$) \textbf{então} Subir$(i)$ \textbf{senão} Descer$(i,n)$.
\end{cuidado}

\begin{exemplo}[{Traços conferidos -- heap $H=[95,60,78,39,28,66,70,33]$}]\leavevmode
\begin{itemize}
  \item \textbf{Aumentar} o nó 6 de 66 para 98: sobe trocando com 78 (pos.\ 3) e com 95 (pos.\ 1): $[98,60,95,39,$\allowbreak$28,78,70,33]$.
  \item \textbf{Remover a raiz}: coloca-se o último (33) na raiz, $[33,60,78,39,28,66,70]$; desce trocando com 78 (maior filho) e depois com 70: $[78,60,70,39,28,66,33]$.
  \item \textbf{Diminuir} a raiz de 95 para 37: desce trocando com 78 e com 70: $[78,60,70,39,$\allowbreak$28,66,37,33]$.
  \item \textbf{Inserir} 73: entra na pos.\ 9 (filho de 39), sobe trocando com 39 e com 60: $[95,73,78,60,$\allowbreak$28,66,70,33,39]$.
\end{itemize}
\end{exemplo}

\subsection{Construção em tempo linear}
\textbf{1ª tentativa} (inserções sucessivas): \textbf{para} $i:=2,\dots,n$ \textbf{faça} subir$(i)$ --- custa $O(n\log n)$ (e $\Theta(n\log n)$ no pior caso, p.ex.\ entrada crescente).

\textbf{2ª tentativa} (\emph{arranjar}): as folhas ($i>\floor{n/2}$) já são heaps; corrige-se de baixo para cima.
\begin{pseudo}{arranjar$(n)$}
\begin{algorithmic}
\Procedure{arranjar}{$n$}
  \For{$i := \floor{n/2}$ \textbf{decr.\ até} $1$} \State \Call{descer}{$i,n$} \EndFor
\EndProcedure
\end{algorithmic}
\end{pseudo}
\emph{Correção:} invariante ``antes de processar $i$, as subárvores de raízes $i+1,\dots,n$ são heaps''. Ao chamar descer$(i,n)$, as duas subárvores dos filhos de $i$ já são heaps, e descer transforma a subárvore de $i$ em heap.

\begin{lema}[Lista de heaps, ex.\ 6]\label{lema:alturas}
Numa árvore binária completa com $n$ nós (folhas com altura 1), no máximo $\ceil{n/2^j}$ nós têm altura $j$.
\end{lema}
\begin{proof}
Numa árvore completa o caminho mais longo a partir de $i$ é o da esquerda: $i\to 2i\to 4i\to\cdots$. Logo o nó $i$ tem altura $\ge j$ sse seu descendente $j-1$ níveis abaixo existe, isto é, sse $i\,2^{j-1}\le n$, ou seja, $i\le\floor{n/2^{j-1}}$. Assim, com $x=n/2^{j-1}$ e $m=\floor{x}$,
\[
\#\{\text{altura}=j\}=\floor{x}-\floor{x/2}=m-\floor{m/2}=\ceil{m/2}\le\ceil{x/2}=\ceil{n/2^j}.\qedhere
\]
\end{proof}

\begin{teo}[Lema 6.1 do Szwarcfiter / slide]
arranjar$(n)$ roda em $O(n)$.
\end{teo}
\begin{proof}
descer$(i,n)$ num nó de altura $j$ faz no máximo $j-1$ trocas. Pelo Lema \ref{lema:alturas}, o total é
\[
\sum_{j=2}^{h}(j-1)\ceil*{\frac{n}{2^j}}\ \le\ \sum_{j\ge2}(j-1)\Bigl(\frac{n}{2^j}+1\Bigr)
\ =\ n\sum_{j\ge2}\frac{j-1}{2^j}+O(\log^2 n)\ =\ n\cdot 1+O(\log^2n)=O(n),
\]
pois $\sum_{j\ge1}\frac{j-1}{2^j}=\sum_{j\ge1}\frac{j}{2^j}-\sum_{j\ge1}\frac1{2^j}=2-1=1$. (O slide faz a mesma conta separando $j$ par e ímpar e usando $\frac{j}{2^{j+1}}\le\frac{j-1}{2^j}$.)
\end{proof}
Intuição: metade dos nós (folhas) não custa nada, $1/4$ custa $\le1$, $1/8$ custa $\le2$, \ldots; só a raiz custa $\log n$.

\begin{exemplo}[Traços do \emph{arranjar} (conferidos)]\leavevmode
\begin{itemize}
  \item Lista de heaps, ex.\ 5: $[18,25,41,34,$\allowbreak$14,10,52,50,48]$ ($n=9$, processa $i=4,3,2,1$).\\
  $i=4$: $34\leftrightarrow50$; $i=3$: $41\leftrightarrow52$; $i=2$: $25\leftrightarrow50$, depois $25\leftrightarrow48$; $i=1$: $18\leftrightarrow52$, depois $18\leftrightarrow41$.\\
  Resultado: $[52,50,41,48,$\allowbreak$14,10,18,34,25]$.
  \item Monitoria (Q4): $[20,15,30,40,10,25,5]$ processa $i=3,2,1$: $i=3$ nada; $i=2$: $15\leftrightarrow40$; $i=1$: $20\leftrightarrow40$. Resultado $[40,20,30,15,10,25,5]$.
  \item Monitoria (Q13): $[7,2,15,9,1,20,13]$: $i=3$: $15\leftrightarrow20$; $i=2$: $2\leftrightarrow9$; $i=1$: $7\leftrightarrow20$, $7\leftrightarrow15$. Resultado $[20,9,15,2,1,7,13]$ (o mesmo que por inserções sucessivas neste caso).
\end{itemize}
\end{exemplo}

\subsection{Heapsort}
\begin{pseudo}{Heapsort (Szwarcfiter, Alg.\ 7.6)}
\begin{algorithmic}
\Procedure{heapsort}{$T,n$}
  \State \Call{arranjar}{$n$} \Comment{$O(n)$}
  \For{$m := n$ \textbf{decrescendo até} $2$}
    \State $T[1]\Leftrightarrow T[m]$ \Comment{o máximo vai para o fim}
    \State \Call{descer}{$1, m-1$} \Comment{$O(\log m)$}
  \EndFor
\EndProcedure
\end{algorithmic}
\end{pseudo}
Complexidade $O(n)+\sum_{m}O(\log m)=\Theta(n\log n)$ em \emph{todos} os casos (mesmo com a entrada ordenada). \textbf{In-place} (só troca no vetor) e \textbf{não estável} (a troca da raiz com o último embaralha iguais). É ótimo entre os algoritmos por comparação. Ex.: $[18,25,41,34,14,10,52,50,48]\to[10,14,18,25,34,41,48,50,52]$.

% ------------------------------------------------------------------
\subsection{Propriedades de heap e suas demonstrações}\label{sec:heapprops}
\begin{prova}[Q3 da AV1: ``saber demonstrar propriedades de heap'']
Quase toda prova de heap usa só três ferramentas: (i) a definição $s_i\le s_{\floor{i/2}}$ + \textbf{transitividade ao longo de um caminho}; (ii) os \textbf{fatos de índice} (pai, filhos, folhas, níveis); (iii) \textbf{indução} (na altura, no nível, ou no número de nós). Para ``provar ou dar contraexemplo'', teste primeiro exemplos pequenos ($n\le 8$).
\end{prova}

\begin{prop}[Ancestral domina]\label{prop:ancestral}
Se $j$ é ancestral de $i$ num heap, então $s_i\le s_j$. Em particular: todo caminho da raiz a uma folha é não crescente, a \textbf{raiz é o máximo} e \textbf{cada subárvore é um heap}.
\end{prop}
\begin{proof}
O caminho de $i$ até $j$ é $i,\floor{i/2},\floor{i/4},\dots,j$; aplicando a definição em cada aresta e a transitividade de $\le$, $s_i\le s_{\floor{i/2}}\le\cdots\le s_j$. A raiz é ancestral de todos. A condição do heap numa subárvore é a mesma condição restrita a ela.
\end{proof}

\begin{prop}[O mínimo está numa folha]
Num heap de máximo com chaves distintas, o menor elemento está numa folha, isto é, numa posição $>\floor{n/2}$.
\end{prop}
\begin{proof}
Se $i$ não é folha, tem filho $2i$ e $s_{2i}<s_i$ (distintas); logo $s_i$ não é o mínimo. (Com repetições: \emph{algum} mínimo está numa folha.)
\end{proof}

\begin{prop}[Onde está o $k$-ésimo maior]
Com chaves distintas, o $k$-ésimo maior elemento está num nível $\le k$, isto é, numa posição $\le 2^k-1$. Em particular o 2º maior é filho da raiz (posição 2 ou 3), e o 3º maior está nas posições $2,\dots,7$.
\end{prop}
\begin{proof}
Pela Prop.~\ref{prop:ancestral}, todos os ancestrais de $x$ são maiores que $x$. Se $x$ é o $k$-ésimo maior, só existem $k-1$ elementos maiores, então $x$ tem no máximo $k-1$ ancestrais, ou seja, está no nível $\le k$. As posições dos níveis $1..k$ são $1,\dots,2^k-1$.
\end{proof}

\begin{prop}[Altura]
Um heap com $n$ elementos tem altura $\floor{\log_2 n}+1$ (em níveis) $=\Theta(\log n)$; por isso subir/descer/inserir/remover são $O(\log n)$.
\end{prop}
\begin{proof}
O último nó, $n$, está no nível $\floor{\log_2n}+1$, e numa árvore completa todos os níveis anteriores estão cheios: $2^{h-1}\le n\le 2^h-1$, logo $h=\floor{\log_2n}+1$.
\end{proof}

\begin{prop}[Relação com ordenação]
Toda lista ordenada de forma não crescente é um heap; a recíproca é falsa (ex.: $[3,1,2]$). Por isso construir um heap é mais fácil que ordenar ($O(n)$ vs.\ $\Omega(n\log n)$) --- e, como heapsort ordena usando construção + $n$ remoções, não existe heap com construção \emph{e} remoção ambas $o(\log n)$ amortizado \emph{por comparação} (senão ordenaríamos em $o(n\log n)$).
\end{prop}

\begin{prop}[Troca de dois elementos -- Lista de heaps, ex.\ 2 e 3]
Seja $S$ um heap e $i<j$.
\begin{enumerate}[label=(\alph*)]
  \item Se $s_i<s_j$, trocar $s_i$ com $s_j$ \textbf{não} produz necessariamente um heap.
  \item Se $s_i>s_j$, trocar $s_i$ com $s_j$ \textbf{não} produz necessariamente um heap.
\end{enumerate}
\end{prop}
\begin{proof}[Contraexemplos (verificados)]
(a) $S=[10,9,1,8,7,0,0,6]$ é heap; $i=3$ ($1$), $j=4$ ($8$). Após a troca, $[10,9,8,1,$\allowbreak$7,0,0,6]$: a posição 8 ($6$) tem pai na posição 4 ($1$) e $6>1$. Intuição: o valor pequeno vai para a posição $j$, cujos filhos podem ser maiores que ele.
(b) $S=[10,9]$, $i=1$, $j=2$: após a troca $[9,10]$, filho maior que o pai. Intuição: o valor pequeno sobe para $i$ e pode ficar abaixo de um descendente maior.
\end{proof}

\begin{prop}[Número de heaps -- Lista de heaps, ex.\ 4]
Seja $H(n)$ o número de permutações de $1,\dots,n$ que são heaps (vale o mesmo para min-heap). Então $H(0)=H(1)=1$ e
\[
H(n)=\binom{n-1}{L}\,H(L)\,H(n-1-L),
\]
onde $L$ é o número de nós da subárvore esquerda da raiz: com $h=\floor{\log_2 n}$ e $u=n-(2^h-1)$ nós no último nível, $L=(2^{h-1}-1)+\min\{u,2^{h-1}\}$. Valores: $1,1,2,3,8,20,80,210,896,3360$ para $n=1..10$ (conferido por força bruta até $n=8$).
\end{prop}
\begin{proof}
A raiz é necessariamente $n$ (máximo). O formato da árvore é fixo, então $L$ é determinado por $n$. Escolhem-se quais $L$ dos $n-1$ valores restantes vão para a esquerda ($\binom{n-1}{L}$ modos); cada lado precisa ser um heap sobre seus valores, e o número de heaps só depende da quantidade de valores ($H(L)$ e $H(n-1-L)$).
\end{proof}

\begin{prop}[Correção de subir e descer]
(a) Se $T$ era heap e a chave de $i$ \emph{aumentou}, após subir$(i)$ $T$ é heap. (b) Se as subárvores dos filhos de $i$ são heaps (em particular se $T$ era heap e a chave de $i$ \emph{diminuiu}), após descer$(i,n)$ a subárvore de $i$ é heap.
\end{prop}
\begin{proof}[Esboço]
(a) Invariante: a única violação possível é entre o nó corrente e seu pai. Ao trocar com o pai $j$ (quando $s_i>s_j$), o antigo $s_j$ desce para $i$ e é $\ge$ aos filhos de $i$ (eram seus descendentes); o novo valor em $j$ é maior que o antigo, então não viola com os filhos de $j$. Repete-se para cima. (b) Invariante: a única violação é entre o nó corrente e seus filhos. Troca-se com o \emph{maior} filho: ele passa a dominar o outro filho e o valor que desceu. A violação desce um nível; termina numa folha ou quando não há violação.
\end{proof}

\begin{chave}
Heap = vetor + ``pai $\ge$ filho''. Subir após aumento, descer após diminuição (maior filho!). Inserir = pôr no fim + subir; remover = último na raiz + descer. Construção linear pelo \emph{arranjar} (de $\floor{n/2}$ até $1$) graças a $\le\ceil{n/2^j}$ nós de altura $j$. Heapsort $\Theta(n\log n)$, in-place, não estável.
\end{chave}

% ================================================================
\section{Heaps avançados: binomiais e de Fibonacci}\label{sec:heapav}
% ================================================================
\textbf{Motivação.} No heap binário (vetor) a \textbf{fusão} (união de duas filas) custa $O(n)$: concatena-se os vetores e reconstrói-se com \emph{arranjar}. Em aplicações com muitas fusões (setores de controle, Dijkstra/Prim em grafos grandes) isso é inviável. O heap binomial faz fusão em $O(\log n)$ pior caso; o de Fibonacci, em $O(1)$ amortizado.

% ------------------------------------------------------------------
\subsection{Árvores binomiais}\label{sec:binomial}
\begin{defi}
$B_0$ é um único nó. Para $k\ge1$, $B_k$ é obtida \textbf{subordinando} (\emph{link}) uma $B_{k-1}$ à outra: a raiz de uma vira filho (mais à direita/novo primeiro filho) da raiz da outra. Num heap, a raiz de \emph{maior} prioridade fica por cima.
\end{defi}

\begin{center}
\begin{tikzpicture}[every node/.style={circle,draw,inner sep=1.3pt,minimum size=4mm}, level distance=7mm]
\node (a) at (0,0) {}; \node[draw=none] at (0,-1.9) {$B_0$};
\node (b) at (1.2,0) {} child {node {}}; \node[draw=none] at (1.2,-1.9) {$B_1$};
\node at (2.9,0) {} [sibling distance=7mm] child {node {} child {node{}}} child {node {}}; \node[draw=none] at (2.9,-1.9) {$B_2$};
\node at (5.6,0) {} [sibling distance=9mm] child {node {} [sibling distance=5mm] child {node{} child{node{}}} child{node{}}} child {node {} child{node{}}} child {node {}}; \node[draw=none] at (5.6,-2.6) {$B_3$};
\end{tikzpicture}
\end{center}

\begin{prop}[Propriedades de $B_k$ -- provas por indução em $k$]\label{prop:binom}\leavevmode
\begin{enumerate}
  \item $B_k$ tem exatamente $2^k$ nós \quad($2\cdot2^{k-1}$).
  \item O grau da raiz (a \textbf{ordem} de $B_k$) é $k$, e é o grau máximo da árvore \quad(o link dá mais um filho à raiz).
  \item A altura de $B_k$ é $k$ (medida em arestas) \quad(a subordinada fica um nível abaixo: $\max\{k-1,\,1+(k-1)\}$).
  \item Há exatamente $\binom kj$ nós na profundidade $j$ \quad($\binom{k-1}{j}+\binom{k-1}{j-1}=\binom kj$; daí o nome).
  \item Os filhos da raiz de $B_k$ são raízes de $B_{k-1},B_{k-2},\dots,B_1,B_0$.
\end{enumerate}
\end{prop}

\begin{cuidado}[Convenção de altura: aula $\times$ livro]
Na aula (18/08) e na monitoria ``Heaps Binomiais'': \textbf{$B_k$ tem altura $k$} (conta arestas). No Szwarcfiter e na monitoria ``Heap-fibonacci-binomial'': \textbf{$B_h$ tem altura $h+1$} (conta níveis, folha tem altura 1, como nos heaps binários do curso). É a mesma árvore! Ao responder exercícios de ``alturas'' (lista, ex.\ 8 e 10), \emph{declare a convenção}. Neste resumo: ordem $=$ grau da raiz $=k$ e altura $=k$ arestas.
\end{cuidado}

% ------------------------------------------------------------------
\subsection{Heap binomial}
\begin{defi}
Um \textbf{heap binomial} é uma floresta $H$ de árvores binomiais tal que (1) cada árvore satisfaz a condição de heap (filho $\le$ pai) e (2) as árvores têm \textbf{ordens distintas} (no máximo uma de cada ordem). As raízes ficam numa lista ordenada por ordem crescente.
\end{defi}

\begin{prop}
(a) Para cada $n$ há um único formato de heap binomial com $n$ nós: $B_k\in H$ sse o bit $k$ de $n$ é 1. (b) Logo $H$ tem no máximo $\floor{\log_2 n}+1$ árvores, cada uma de altura $\le\floor{\log_2n}$.
\end{prop}
\begin{proof}
$B_k$ contribui $2^k$ nós e as ordens são distintas, então $n=\sum_{B_k\in H}2^k$ --- é a representação binária de $n$, que é única. Um número $n$ tem $\floor{\log_2 n}+1$ bits.
\end{proof}
\noindent Ex.: $13=1101_2\Rightarrow H=\{B_0,B_2,B_3\}$; $26=11010_2\Rightarrow\{B_1,B_3,B_4\}$.

\subsubsection{Operações}
\begin{pseudo}{Binomial-Link e fusão (Szwarcfiter Alg.\ 9 / monitoria)}
\begin{algorithmic}
\Procedure{Binomial-Link}{$y,z$} \Comment{$y$ e $z$ raízes de $B_{k-1}$; $z$ fica por cima}
  \State $p[y]\gets z$;\quad $sibling[y]\gets child[z]$;\quad $child[z]\gets y$;\quad $degree[z]\gets degree[z]+1$
\EndProcedure
\Procedure{Fusão}{$H',H''$}
  \State $H := H'\cup H''$ \Comment{intercala as listas de raízes por ordem crescente: $O(\log n)$}
  \While{$H$ contém duas árvores $B,B'$ de mesma ordem} \Comment{a menor ordem repetida primeiro}
    \If{$B.raiz.chave \ge B'.raiz.chave$} \State subordinar $B'$ a $B$
    \Else \State subordinar $B$ a $B'$
    \EndIf
  \EndWhile
\EndProcedure
\end{algorithmic}
\end{pseudo}
A fusão é exatamente \textbf{uma soma binária com ``vai um''}: duas $B_k$ viram uma $B_{k+1}$ (o ``vai um''). Como há $O(\log n)$ posições e no máximo 3 árvores de mesma ordem ao mesmo tempo (duas originais + um carry), a fusão custa $O(\log n)$ no pior caso. (Na implementação do Cormen, percorre-se a lista com ponteiros \emph{prev-x}, \emph{x}, \emph{next-x} e há 4 casos: ordens diferentes; três iguais seguidas; dois iguais com $x$ vencendo; dois iguais com \emph{next-x} vencendo.)

\begin{exemplo}[Fusão -- Szwarcfiter Fig.\ 9.10 e exercício da monitoria]
$H_1=\{B_0,B_2,B_4\}$ ($21=10101_2$) e $H_2=\{B_0,B_1\}$ ($3=11_2$). Links na ordem: $B_0+B_0\to B_1$; esse $B_1+B_1$ (de $H_2$) $\to B_2$; esse $B_2+B_2$ (de $H_1$) $\to B_3$. Resultado $\{B_3,B_4\}$, isto é, $24=11000_2=10101_2+11_2$.
\end{exemplo}

\begin{itemize}
  \item \textbf{Inserção:} fusão com um heap $\{B_0\}$ contendo só o novo nó. Pior caso $O(\log n)$ (quando $n=2^k-1$, todos os bits são 1 e há $k$ links). \emph{Tempo esperado $O(1)$} (Lema 9.3 do Szwarcfiter: se $j$ é a posição do primeiro 0 de $n$ em binário, lida da direita para a esquerda, a inserção faz $j$ ``fusões'' no sentido do livro, das quais $j-1$ são links, um por 1 final de $n$; com bits equiprováveis, $E[j]=\sum_{j\ge1}j/2^j=2$, isto é, 2 passos esperados e 1 link esperado) e \emph{amortizado $O(1)$} (dá-se 1 crédito a cada árvore; cada link é pago pelo crédito da árvore subordinada, que some --- é o contador binário).
  \item \textbf{Seleção do máximo:} o máximo é uma das $\le\floor{\log n}+1$ raízes: $O(\log n)$, ou $O(1)$ mantendo um ponteiro para a maior raiz (atualizado nas operações).
  \item \textbf{Remoção do máximo:} ache a raiz $r$ máxima, numa árvore $B_k$; retire $B_k$ de $H$ ($H'=H-B_k$); retirando $r$, seus filhos $B_{k-1},\dots,B_0$ formam um heap binomial $H''$; faça a fusão $H'\cup H''$. Total $O(\log n)$.
  \item \textbf{Alteração de prioridade:} \emph{aumentar} $\Rightarrow$ sobe dentro da própria árvore trocando com o pai, como no heap binário: $O(\log n)$ (altura $\le\log n$). \emph{Diminuir} $\Rightarrow$ remover o nó (ver abaixo) e reinseri-lo com a nova prioridade: $O(\log n)$. (Descer trocando com o maior filho também funciona, mas cada nível tem até $O(\log n)$ filhos para comparar: $O(\log^2 n)$.)
  \item \textbf{Remoção de um nó qualquer:} aumente sua prioridade para $+\infty$ (sobe até a raiz) e remova o máximo: $O(\log n)$.
  \item \textbf{Construção:} $n$ inserções custam $O(n)$ no total (contador binário: $\sum_k n/2^k\le 2n$). É a resposta do ex.\ 7 da lista.
\end{itemize}

% ------------------------------------------------------------------
\subsection{Heap de Fibonacci}\label{sec:fib}
\textbf{Ideia: preguiça.} O heap binomial é ``estrito'': toda inserção/fusão já resolve colisões de ordem. O de Fibonacci \emph{adia} toda a arrumação para a remoção do máximo.

\paragraph{Estrutura.} Floresta de árvores com condição de heap; as raízes numa \textbf{lista circular duplamente encadeada}, e os filhos de cada nó também. Ponteiro para a raiz máxima. Cada nó guarda: chave, pai, um filho, irmãos esquerdo/direito, grau, e um bit \textbf{marca}. Listas circulares duplas permitem inserir/remover um nó ou concatenar duas listas em $O(1)$.

\paragraph{Operações} (heap de máximo, como no curso; ``aumentar prioridade'' = \emph{decrease-key} do Cormen, que usa heap de mínimo).
\begin{itemize}
  \item \textbf{Inserção:} cria uma árvore de um nó e a coloca na lista de raízes; atualiza o ponteiro do máximo. $O(1)$ real.
  \item \textbf{Fusão:} concatena as duas listas de raízes; máximo $=$ maior dos dois. $O(1)$ real.
  \item \textbf{Seleção:} ponteiro. $O(1)$.
  \item \textbf{Remoção do máximo:} (1) remove a raiz máxima e promove todos os seus filhos a raízes (desmarcados); (2) \textbf{consolida}: percorre as raízes com um vetor $A[0..D]$ indexado pelo grau ($D=O(\log n)$); sempre que duas raízes têm o mesmo grau, subordina a de menor prioridade à de maior (o grau sobe de 1) e repete; (3) reconstrói a lista de raízes a partir de $A$ e acha o novo máximo. Real: $O(D+t)$ com $t=$ número de raízes; \textbf{amortizado $O(\log n)$}.
  \item \textbf{Aumento de prioridade de $x$:} aumenta a chave; se não ficou maior que o pai, fim. Senão: \textbf{corta} $x$ (com sua subárvore) e o põe na lista de raízes, desmarcado; então, subindo: se o pai $y$ está \emph{desmarcado}, marque-o (se não for raiz) e pare; se $y$ já estava \emph{marcado} (perdeu um filho antes), corte $y$ também, desmarque-o, leve-o à lista de raízes e continue com o pai de $y$ --- \textbf{cortes em cascata}. \textbf{Amortizado $O(1)$}.
  \item \textbf{Remoção de nó qualquer:} aumenta para $+\infty$ e remove o máximo: $O(\log n)$\am.
  \item \textbf{Diminuição de prioridade:} não é a operação ``barata''; faz-se removendo e reinserindo (ou cortando os filhos que passarem a violar): $O(\log n)$\am. Na tabela da aula (convenção de mínimo), isso aparece como ``incrementar $O(\log n)$\am, decrementar $O(1)$\am''.
\end{itemize}

\paragraph{Regras das marcas.} Todos começam desmarcados; um nó é marcado quando \emph{perde o primeiro filho} por corte; ao perder o segundo, ele próprio é cortado. \textbf{Raízes nunca estão marcadas} (e podem perder vários filhos); ao ser subordinado a outro nó, o nó fica desmarcado.

\begin{exemplo}[Cortes em cascata -- exercício da monitoria de 26/08]
Nó interno $X$ com pai $Y$ (marcado) e avô $Z$ (marcado); a prioridade de $X$ aumenta violando a ordem. (1) $X$ é cortado e vai para a lista de raízes (desmarcado). (2) $Y$ perdeu seu segundo filho: é cortado, desmarcado, vai para a lista de raízes. (3) Idem para $Z$. (4) Seja $W$ o pai de $Z$: se $W$ é raiz, nada muda; se não é raiz e está desmarcado, passa a \textbf{marcado} e a cascata para; se estava marcado, a cascata continua. Nós movidos para as raízes: $X,Y,Z$ (e possivelmente mais); flags: $X,Y,Z$ desmarcados, $W$ marcado (se não raiz).
\end{exemplo}

\subsubsection{Por que funciona: grau máximo $O(\log n)$}
\begin{lema}
Seja $x$ um nó e $y_1,\dots,y_k$ seus filhos na ordem em que foram ligados a $x$. Então $\operatorname{grau}(y_i)\ge i-2$.
\end{lema}
\begin{proof}
Quando $y_i$ foi ligado, $x$ já tinha $y_1,\dots,y_{i-1}$, logo grau $\ge i-1$; só se ligam nós de mesmo grau, então $y_i$ tinha grau $\ge i-1$. Desde então $y_i$ perdeu no máximo um filho (senão teria sido cortado). Logo $\operatorname{grau}(y_i)\ge i-2$.
\end{proof}
\begin{teo}
Um nó de grau $k$ tem subárvore com pelo menos $F_{k+2}\ge\varphi^k$ nós ($F$ = Fibonacci, $\varphi=\frac{1+\sqrt5}2$). Logo o grau máximo é $D(n)\le\log_\varphi n=O(\log n)$, e um heap de Fibonacci tem $O(\log n)$ árvores após a consolidação.
\end{teo}
\begin{proof}[Esboço]
Seja $s_k$ o tamanho mínimo de uma subárvore de raiz com grau $k$. Pelo lema, $s_k\ge 2+\sum_{i=2}^{k}s_{i-2}$, e por indução $s_k\ge F_{k+2}$ (pois $F_{k+2}=1+\sum_{i=0}^{k}F_i$). É daí que vem o nome: as menores árvores de posto $k$ ($F_0,\dots,F_5$ desenhadas em aula) têm $1,2,3,5,8,13$ nós.
\end{proof}
\noindent Sem marcas/cascata, poderíamos cortar todos os netos de uma raiz: uma árvore de posto $k$ com só $k+1$ nós; o posto deixaria de ser $O(\log n)$ e a consolidação ficaria cara (resposta ao ``Exercício 3'' da monitoria).

\subsubsection{Análise amortizada}
Potencial (Cormen, cap.~19): $\Phi(H)=t(H)+2\,m(H)$, com $t=$ nº de raízes e $m=$ nº de nós marcados.
\begin{itemize}
  \item Inserção: real $O(1)$, $\Delta\Phi=+1$ $\Rightarrow O(1)$. Fusão: $\Delta\Phi=0$ $\Rightarrow O(1)$.
  \item Remoção do máximo: real $O(D+t)$; depois há $\le D+1$ raízes, então $\Delta\Phi\le (D+1)-t$ $\Rightarrow$ amortizado $O(D)=O(\log n)$. \emph{O custo de consolidar as $t$ raízes acumuladas foi ``pago'' pelas inserções/fusões preguiçosas.}
  \item Aumento de prioridade com $c$ cortes em cascata: real $O(c)$; $t$ cresce $c$, $m$ cai $\ge c-2$ $\Rightarrow\Delta\Phi\le c+2(2-c)=4-c$ $\Rightarrow$ amortizado $O(1)$. (Visão do banqueiro do Szwarcfiter: o corte de um nó marcado é debitado à operação que o marcou.)
\end{itemize}
\begin{teo}[Aula de 21/08]
A partir de um heap de Fibonacci vazio, qualquer sequência de $a_1$ inserções, $a_2$ remoções, $a_3$ aumentos de prioridade e $a_4$ uniões custa $O(a_1+a_3+a_4+a_2\log n)$.
\end{teo}

\begin{cuidado}
\begin{itemize}
  \item Os $O(1)$ do Fibonacci são \textbf{amortizados}: \emph{uma} remoção do máximo após muitas inserções/fusões pode custar $\Theta(n)$ (consolida tudo de uma vez). Se o requisito é previsibilidade de cada operação, Fibonacci é ruim.
  \item Constantes grandes e muita memória por nó (4 ponteiros + grau + marca): na prática só compensa com MUITAS alterações de prioridade (Dijkstra/Prim em grafos densos).
  \item Uma árvore de um heap de Fibonacci \emph{pode} ter altura grande (até $\Theta(n)$ em sequências adversárias; Cormen, ex.~19.4-1): o que é logarítmico é o \textbf{grau}, não a altura.
\end{itemize}
\end{cuidado}

% ------------------------------------------------------------------
\subsection{Comparação e escolha da estrutura (Q4 e Teste 1)}\label{sec:escolhaheap}
\begin{center}
\small
\begin{tabular}{@{}l>{\centering\arraybackslash}p{3.3cm}>{\centering\arraybackslash}p{4.2cm}>{\centering\arraybackslash}p{3.6cm}@{}}
\toprule
& Heap binário & Heap binomial & Heap de Fibonacci\\
\midrule
Seleção & $O(1)$ & $O(\log n)$ / $O(1)$ c/ ponteiro & $O(1)$\\
Inserção & $O(\log n)$ & $O(\log n)$ pior; $O(1)$ esp./amort. & $O(1)$\am\\
Remoção do máx. & $O(\log n)$ & $O(\log n)$ & $O(\log n)$\am\\
Aumentar prioridade & $O(\log n)$ & $O(\log n)$ & $O(1)$\am\\
Fusão & $O(n)$ & $O(\log n)$ \textbf{pior caso} & $O(1)$\am\\
Construção & $O(n)$ & $O(n)$ & $O(n)$\\
Memória & vetor, sem ponteiros & 3 ponteiros + grau/nó & 4 ponteiros + grau + marca/nó\\
Garantia & pior caso & pior caso & amortizada\\
Simplicidade & muito simples & média & complexa\\
\bottomrule
\end{tabular}
\end{center}

\begin{prova}[Roteiro de escolha]
\begin{enumerate}
  \item Precisa de \textbf{fusão}? Não $\Rightarrow$ \textbf{heap binário} (mais simples, menos memória, ótimo em cache).
  \item Fusão frequente e \textbf{cada operação precisa ser previsível} (limite de tempo por operação, sistema de tempo real) $\Rightarrow$ \textbf{heap binomial} (fusão $O(\log n)$ \emph{pior caso}).
  \item Muitíssimos \textbf{aumentos de prioridade} (ou fusões) e o que importa é o \textbf{custo total} da execução $\Rightarrow$ \textbf{heap de Fibonacci}.
  \item Restrição de \textbf{memória} forte $\Rightarrow$ binário (as outras têm overhead de ponteiros por nó).
\end{enumerate}
\end{prova}

\paragraph{Teste 1 -- respostas esperadas.}
\begin{itemize}
  \item \textbf{Q1 (aeroporto):} fusões muito frequentes com filas de centenas de milhares de aeronaves; alterações raras; exigência de tempo \emph{previsível por operação individual}. $\Rightarrow$ \textbf{Heap binomial}: fusão em $O(\log n)$ no pior caso de \emph{cada} execução, crescendo só logaritmicamente. Binário: fusão $O(n)$ (reconstrução) --- milhares de fusões com $10^5$ elementos estouram a janela. Fibonacci: fusão $O(1)$ só \emph{amortizada}; o custo acumulado aparece de uma vez numa remoção (consolidação $\Theta(n)$ possível) --- imprevisível; e sua vantagem (aumento de prioridade $O(1)$) não é usada, pois alterações são raras.
  \item \textbf{Q2 (central médica, 40 MB):} $2\cdot10^6$ chamadas $\times$ 16 bytes $=32\,000\,000$ bytes $\approx 30{,}5$ MB $<41\,943\,040$. Heap binomial precisaria de $\ge3$ ponteiros (pai, filho, irmão $=24$ bytes) $+$ grau: $\ge 40$ bytes/nó $\Rightarrow\ge 80\cdot10^6$ bytes; Fibonacci ($\ge4$ ponteiros $+$ grau $+$ marca) ainda mais. Não há fusão, alterações são poucas, inserção/remoção $O(\log n)$ bastam. $\Rightarrow$ \textbf{Heap binário em vetor} (único que cabe, e o mais simples).
  \item \textbf{Q3 (supercomputador):} $5\cdot10^6$ inserções, $5\cdot10^6$ remoções, $80\cdot10^6$ aumentos de urgência, poucas uniões; importa o custo total. Fibonacci: $\approx 5\cdot10^6+80\cdot10^6+5\cdot10^6\log_2(5\cdot10^6)\approx 2\cdot10^8$ passos; binário/binomial: $\approx 90\cdot10^6\cdot 22\approx 2\cdot10^9$. $\Rightarrow$ \textbf{Heap de Fibonacci} (aumento $O(1)$\am\ sem remover/reinserir). Binário/binomial pagam $O(\log n)$ em cada um dos 80 milhões de aumentos.
\end{itemize}

% ================================================================
\section{Conjuntos disjuntos: Union-Find}\label{sec:uf}
% ================================================================

\subsection{O problema}
Manter uma partição de um universo de $n$ elementos em conjuntos disjuntos, cada um identificado por um \textbf{representante}, com:
\begin{itemize}
  \item \MakeSet$(x)$: cria o conjunto $\{x\}$ (Szwarcfiter: \emph{gerar});
  \item \Find$(x)$: devolve o representante do conjunto de $x$ (\emph{buscar}); ``$x$ e $y$ estão juntos?'' $\iff\Find(x)=\Find(y)$;
  \item \Union$(x,y)$: une os conjuntos de $x$ e $y$ (\emph{fundir}).
\end{itemize}
Conjuntos \textbf{nunca se separam} (não há ``split''). Uma sequência típica: $n$ \MakeSet{} e $m$ \Find/\Union{} intercalados; no máximo $n-1$ uniões efetivas. O número de conjuntos $=$ número de raízes, que nunca aumenta.
\textbf{Aplicações:} componentes conexas dinâmicas, algoritmo de Kruskal (AGM), equivalência de símbolos em compiladores, grupos em redes sociais, percolação.

\subsection{Implementações}
\subsubsection*{(a) Vetor de representantes (``quick-find'')}
$rep[x]$ guarda o representante. \Find$=O(1)$; \Union{} precisa trocar $rep$ de todos os elementos de um conjunto: $O(n)$ por união.

\subsubsection*{(b) Listas encadeadas com ponteiro direto ao representante}
Cada conjunto é uma lista; um objeto-conjunto guarda \emph{head}, \emph{tail} e \emph{size}; todo nó tem ponteiro para o objeto-conjunto (o representante).
\begin{itemize}
  \item \MakeSet: $O(1)$. \Find: $O(1)$ (segue um ponteiro).
  \item \Union: concatena as listas em $O(1)$ (tail de uma $\to$ head da outra) \emph{mas} precisa atualizar o ponteiro ao representante de todos os nós da lista anexada.
  \item Sem cuidado: $\Theta(n)$ por união, $\Theta(n^2)$ numa sequência ruim (sempre anexar a lista grande à unitária).
  \item \textbf{União ponderada} (anexa a \emph{menor} à maior): cada vez que o ponteiro de um elemento $x$ é atualizado, o conjunto de $x$ pelo menos \textbf{dobra}; como tamanhos $\le n$, $x$ é atualizado $\le\log_2 n$ vezes. Total de $n$ \MakeSet{} e $m$ operações: $O(m+n\log n)$, isto é, \Union{} em $O(\log n)$ \textbf{amortizado}.
  \item Vantagem exclusiva: \textbf{listar os membros} de um conjunto em $O(k)$ ($k$ = tamanho do conjunto), percorrendo a lista.
\end{itemize}

\subsubsection*{(c) Floresta (``quick-union'')}
Cada conjunto é uma árvore; cada nó aponta para o pai ($pai[x]$); a raiz é o representante ($pai[r]=r$ ou $0$). \MakeSet: $pai[x]:=x$. \Find: sobe até a raiz, $O(h)$. \Union: acha as duas raízes e faz uma apontar para a outra, $O(1)$ \emph{dadas as raízes}.
Problema: sem critério, uma sequência de uniões gera uma \textbf{cadeia de altura $n-1$} (Find $\Theta(n)$).

\subsubsection*{(d) União por tamanho ou por rank}
Liga-se a raiz da árvore \emph{menor} (em tamanho, ou em \emph{rank}) à raiz da maior. \textbf{Rank} (aula de 25/08):
\begin{itemize}
  \item todo nó começa com rank 0;
  \item só o rank de raízes é atualizado: na união de raízes de ranks \emph{iguais} $r$, a nova raiz passa a ter rank $r+1$; com ranks diferentes, nada muda;
  \item o rank nunca decresce; quando um nó deixa de ser raiz, seu rank nunca mais muda;
  \item sem compressão de caminhos, rank $=$ altura; com compressão, rank é só uma cota superior da altura.
\end{itemize}

\begin{teo}[Altura logarítmica]\label{teo:ufaltura}
Com união por rank (ou por tamanho), toda árvore de rank/altura $h$ tem pelo menos $2^h$ nós. Logo $h\le\floor{\log_2 n}$ e \Find{} custa $O(\log n)$.
\end{teo}
\begin{proof}
Indução no número de operações. Uma árvore de um nó tem $h=0$ e $2^0=1$ nó. A altura (rank) de uma árvore só aumenta quando se unem duas de mesmo rank $h$: pela hipótese cada uma tem $\ge2^h$ nós, e a nova tem $\ge 2\cdot2^h=2^{h+1}$ nós e rank $h+1$. Ligando a menor sob a maior, a altura não aumenta. Como $2^h\le n$, $h\le\log_2 n$.
(Szwarcfiter, Teo.~8.2, com união por tamanho e altura contada em níveis: $2^{h(T)-1}\le|T|$, usando $h(T)=\max\{1+h(T_v),h(T_w)\}$ com $|T_v|\le|T_w|$.)
\end{proof}

\begin{cor}
(i) O número de nós com rank exatamente $k$ é $\le n/2^k$ (cada um é raiz, em algum momento, de uma subárvore com $\ge2^k$ nós, e essas subárvores são disjuntas). (ii) Todo rank é $\le\floor{\log_2n}$. (iii) $\rank(x)<\rank(pai[x])$ para todo $x$ não raiz.
\end{cor}

\subsubsection*{(e) Compressão de caminhos}
Durante \Find$(x)$, depois de achar a raiz, faz-se \emph{todo nó do caminho} apontar diretamente para a raiz. As próximas buscas desses nós custam $O(1)$. Variantes (Szwarcfiter): \emph{separação} de caminhos (cada nó passa a apontar para o avô) e \emph{divisão} (nós alternados apontam para o avô) --- mesma complexidade assintótica, uma só passada.

\begin{pseudo}{Floresta com união por rank e compressão de caminhos}
\begin{algorithmic}
\Procedure{Make-Set}{$x$} \State $pai[x]:=x$;\ $rank[x]:=0$ \EndProcedure
\Function{Find}{$x$}
  \If{$pai[x]\ne x$} \State $pai[x]:=$ \Call{Find}{$pai[x]$} \Comment{compressão}
  \EndIf
  \State \Return $pai[x]$
\EndFunction
\Procedure{Union}{$x,y$}
  \State $r:=$ \Call{Find}{$x$};\ $s:=$ \Call{Find}{$y$}
  \If{$r=s$} \Return \EndIf
  \If{$rank[r]<rank[s]$} \State $pai[r]:=s$
  \ElsIf{$rank[r]>rank[s]$} \State $pai[s]:=r$
  \Else \State $pai[s]:=r$;\ $rank[r]:=rank[r]+1$
  \EndIf
\EndProcedure
\end{algorithmic}
\end{pseudo}
Versão iterativa da compressão (Szwarcfiter, \emph{comp-buscar}): uma passada acha a raiz $y$; uma segunda passada percorre de novo o caminho fazendo $pai[\cdot]:=y$.

\subsection{Rank + compressão: $O(m\lgs n)$ (demonstração da aula de 25/08)}
$\lgs n$ = número de vezes que se aplica $\log_2$ até chegar a $\le1$. Defina $F(0)=1$, $F(k)=2^{F(k-1)}$: $F=1,2,4,16,65536,2^{65536}$. Então $\lgs n$ é o menor $k$ com $F(k)\ge n$ ($\lgs n\le5$ para qualquer $n$ prático).

\textbf{Grupos de ranks:} o grupo $B\ge1$ contém os ranks em $(F(B-1),F(B)]$, e o grupo 0 os ranks $\le F(0)=1$: grupo 0 $=\{0,1\}$, 1 $=\{2\}$, 2 $=\{3,4\}$, 3 $=\{5,\dots,16\}$, 4 $=\{17,\dots,65536\}$.
\begin{itemize}
  \item \textbf{Propriedade A:} como ranks $\le\log n$, há no máximo $\lgs n$ grupos (ranks até $\log n$ caem no grupo $\le\lgs n-1$).
  \item \textbf{Propriedade B:} o número de nós no grupo $B$ é $\le\sum_{r=F(B-1)+1}^{F(B)}\frac{n}{2^r}\le\frac{n}{2^{F(B-1)}}=\frac{n}{F(B)}$ (para $B\ge1$; o grupo 0 tem trivialmente $\le n=n/F(0)$ nós).
\end{itemize}
Num \Find, cobre-se cada nó $x$ do caminho: se $x$ é raiz, filho da raiz, ou se $pai[x]$ está num \emph{grupo diferente} do de $x$, cobra-se do \Find{} --- como os ranks crescem ao longo do caminho, há $O(\lgs n)$ mudanças de grupo por \Find: total $O(m\lgs n)$. Caso contrário ($pai[x]$ no mesmo grupo), cobra-se de $x$: a compressão dá a $x$ um novo pai de rank \emph{estritamente maior}; um nó do grupo $B$ só pode receber isso $\le F(B)$ vezes antes que seu pai vá para outro grupo (e aí nunca mais é cobrado). Total por grupo: $\le\frac{n}{F(B)}\cdot F(B)=n$; somando os $\lgs n$ grupos, $O(n\lgs n)$. \textbf{Total: $O(m\lgs n)$} (com $m\ge n$).

\begin{cuidado}[Deslocamento na lista de grupos das notas]
Nas notas de 25/08 a lista explícita ``grupo 0: $\{0\}$, grupo 1: $\{1\}$, grupo 2: $\{2\}$, grupo 3: $\{3,4\}$, grupo 4: $\{5,\dots,16\}$'' está deslocada de uma posição em relação à fórmula $(F(B-1),F(B)]$ escrita logo acima, que dá grupo 1 $=(1,2]=\{2\}$, grupo 2 $=(2,4]=\{3,4\}$ etc. Nada muda na análise: continua havendo $O(\lgs n)$ grupos e a Propriedade B continua valendo. Na prova, use a fórmula e seja consistente. As notas escrevem $m$ tanto para o número de elementos quanto para o de operações; aqui usamos $n$ para elementos e $m$ para operações.
\end{cuidado}

\begin{obs}
O limite justo (Tarjan) é $O(m\,\alpha(m,n))$, com $\alpha$ a inversa da função de Ackermann ($\alpha\le4$ na prática; Szwarcfiter §8.4: $O(n+m\,\alpha(m+n,n))$). Para a prova, o resultado demonstrado em aula é o $O(m\lgs n)$; ambos são ``quase lineares''.
\end{obs}

% ------------------------------------------------------------------
\subsection{Qual implementação usar? (Q2 da AV1)}\label{sec:ufescolha}
\begin{table}[h]
\centering\small
\caption{Implementações de Union-Find e quando usar cada uma.}\label{tab:uf}
\begin{tabularx}{\linewidth}{@{}p{3.3cm}cccX@{}}
\toprule
Implementação & \MakeSet & \Find & \Union & Quando é a melhor escolha\\
\midrule
Vetor de representantes & $O(1)$ & $O(1)$ & $O(n)$ & Quase só consultas e pouquíssimas uniões, $n$ pequeno.\\
Listas encadeadas + ponteiro ao representante + união ponderada & $O(1)$ & $O(1)$ \textbf{pior caso} & $O(\log n)$\am & Consultas dominam e cada uma precisa ser $O(1)$ \emph{garantido}; é preciso \textbf{listar os membros} de um grupo em $O(k)$; memória extra aceitável (head/tail/size + ponteiros).\\
Floresta ingênua & $O(1)$ & $O(n)$ & $O(1)$\,+\,Finds & Nunca (só didática).\\
Floresta + união por rank/tamanho (sem compressão) & $O(1)$ & $O(\log n)$ & $O(\log n)$ & Quando é preciso \textbf{desfazer uniões} (rollback/``split'' offline): sem compressão, cada união muda só um ponteiro (e talvez um rank), que se empilha e se desfaz em $O(1)$.\\
Floresta + rank + compressão de caminhos & $O(1)$ & \multicolumn{2}{c}{$O(\lgs n)$\am\ ($\alpha$)} & Caso geral, muitas operações misturadas, só \emph{Find} e \emph{Union} (Kruskal, conectividade, redes sociais): praticamente $O(1)$ por operação, memória mínima (um vetor $pai$ + $rank$).\\
\bottomrule
\end{tabularx}
\end{table}

\begin{prova}[Cenários típicos (documento ``Conjuntos Disjuntos'' do Classroom)]
\begin{itemize}
  \item \textbf{``Listar os membros do grupo de $U$''}: a floresta só tem ponteiros filho$\to$pai; listar exigiria testar $\Find(i)=\Find(U)$ para todos os $N$ elementos ($O(N)$). Com \textbf{listas encadeadas}: percorre-se a lista, $O(k)$; \emph{MesmoGrupo} em $O(1)$; \emph{Juntar} em $O(\log n)$\am{} (menor na maior).
  \item \textbf{Remoção de arestas / conectividade após destruições} (``A queda do império''): Union-Find não desfaz uniões (e a compressão destrói a topologia). Solução \emph{offline}: leia todas as consultas, construa o estado final (sem as arestas que serão destruídas) e processe \textbf{de trás para frente}, transformando cada destruição numa \Union; guarde as respostas e imprima invertido.
  \item \textbf{Operações online com desfazer (snapshot/rollback)}: \textbf{sem compressão}, só união por rank/tamanho (altura $O(\log n)$, \Find{} $O(\log n)$); cada \Union{} empilha o ponteiro/rank alterado; desfazer $=$ desempilhar e restaurar em $O(1)$.
  \item \textbf{Kruskal}: ordena arestas; para cada aresta $(u,v)$, se $\Find(u)\ne\Find(v)$ inclui e faz \Union. Floresta com rank + compressão.
\end{itemize}
\end{prova}

\begin{prova}[Teste 1, Q4 -- plataforma com $10^7$ usuários]
$10^5$ uniões e $10^8$ consultas ``mesmo grupo?''; grupos nunca se separam; não importa o histórico.
\textbf{(a)} Estrutura: \textbf{Union-Find}. Entre ``árvore apontando até o representante'' e ``lista com ponteiro direto ao representante'': as consultas dominam ($10^8\gg10^5$) e precisam ser rápidas; na \textbf{lista encadeada} cada \Find{} é $O(1)$ \emph{no pior caso} (um acesso), e o custo fica nas raras uniões ($O(\log n)$\am{} com união ponderada: no total $\le n\log n\approx 10^7\cdot 23$ atualizações, diluídas). \textbf{Resposta recomendada: a lista encadeada}, porque com razão consultas:uniões de $10^8:10^5$ o que importa é o custo de cada \Find, e a lista o garante em $O(1)$ \emph{no pior caso}. A floresta com rank + compressão também é defensável ($O(\alpha)$\am{} por operação e menos memória), desde que se explique que a garantia é amortizada; o que \emph{não} se deve propor é floresta sem balanceamento.
\textbf{(b)} Estratégia de união: \textbf{anexar o menor ao maior} (união por tamanho na lista; por rank/tamanho na árvore). Na lista, garante que cada elemento muda de representante $\le\log_2 n$ vezes (o tamanho do seu grupo dobra a cada mudança). Na árvore, garante altura $\le\log_2 n$ (Teo.~\ref{teo:ufaltura}), evitando cadeias de altura $\Theta(n)$ que tornariam cada consulta $\Theta(n)$.
\end{prova}

\subsection{Exercícios resolvidos (monitoria de revisão e lista)}
\begin{exemplo}[Revisão, Q2]\leavevmode
Partindo de $\{A\},\dots,\{F\}$ e as uniões \Union$(A,B)$, \Union$(C,D)$, \Union$(A,C)$, \Union$(E,F)$, com união por rank e, em empate, a primeira raiz vence.
\Union$(A,B)$: ranks $0=0\Rightarrow$ raiz $A$, $\rank(A)=1$. \Union$(C,D)$: raiz $C$, $\rank(C)=1$. \Union$(A,C)$: $1=1\Rightarrow$ raiz $A$, $\rank(A)=2$ (árvore $A\to\{B,C\}$, $C\to D$). \Union$(E,F)$: raiz $E$, $\rank(E)=1$. Representantes finais: $A$ (de $\{A,B,C,D\}$) e $E$ (de $\{E,F\}$): \textbf{2 conjuntos}.
\end{exemplo}
\begin{exemplo}[Revisão, Q3 e Q4]
Cadeia $pai[B]=A$, $pai[C]=B$, $pai[D]=C$, $pai[E]=D$. \Find$(E)$ com compressão devolve $A$ e deixa $B,C,D,E$ todos apontando para $A$ (árvore de altura 1). Vantagem: consultas futuras a esses nós custam $O(1)$.
Se $\rank(A)=\rank(D)=1$ e \Union$(A,D)$: empate, a nova raiz é $A$ (pela convenção) e $\rank(A)=2$. O rank só sobe no empate porque, ligando a menor sob a maior, a altura da maior não muda; só com alturas iguais a nova árvore fica 1 mais alta.
\end{exemplo}
\begin{exemplo}[Revisão, Q5 -- por que usar as duas técnicas juntas?]
Só rank, sem compressão: \Find{} fica $\Theta(\log n)$ (há sequências que atingem essa altura). Só compressão, sem rank: uma única árvore pode ficar alta (cadeia) antes de ser comprimida; custo total $O(m\log n)$ (Tarjan--van Leeuwen), e \Find{} individual pode custar $\Theta(n)$. Juntas: $O(m\lgs n)$ (ou $O(m\,\alpha)$) no total --- praticamente constante por operação.
\end{exemplo}
\begin{exemplo}[Lista UF/AVL, ex.\ 1 -- a cota $\floor{\log n}$ é justa]
A prova é o Teo.~\ref{teo:ufaltura}. Sequência que atinge altura $\log_2 n$ ($n=2^k$): una os elementos aos pares (altura 1), depois as árvores aos pares (altura 2), etc. Após $k$ rodadas há uma árvore de rank/altura $k=\log_2n$ (é uma árvore binomial $B_k$!).
\end{exemplo}
\begin{exemplo}[Lista UF/AVL, ex.\ 2 -- só compressão]
Árvore alta antes da compressão: \Union$(1,2)$, \Union$(1,3)$, \dots, sempre pendurando a raiz da árvore grande sob o novo elemento unitário ($pai[\text{raiz}]:=i$): forma uma cadeia $1\to2\to\cdots\to n$ de altura $n-1$, sem nenhum \Find{} caro no caminho. Um \Find$(1)$ custa $\Theta(n)$ e achata tudo. Resultado conhecido: com compressão e união arbitrária, $m$ operações custam $O(m\log n)$ (na verdade $\Theta(m\log_{1+m/n}n)$); a ideia é que cada compressão que passa por um nó ``leve'' (cuja subárvore é menos da metade da do pai) só pode acontecer $O(\log n)$ vezes por nó.
\end{exemplo}

\begin{chave}
Q2 em uma frase: \textbf{listas encadeadas} dão \Find{} $O(1)$ pior caso, \Union{} $O(\log n)$\am\ com ``menor na maior'' e listagem de membros em $O(k)$; \textbf{florestas} com \textbf{rank + compressão} dão $O(\lgs n)$ (ou $\alpha$) amortizado para tudo com memória mínima; \textbf{sem compressão} se for preciso desfazer uniões. Sempre \textbf{menor no maior}, senão árvores/listas degeneram.
\end{chave}

% ================================================================
\section{Árvores binárias de busca e árvores AVL}\label{sec:avl}
% ================================================================

\subsection{Árvore binária de busca (ABB)}
\textbf{Motivação:} lista sequencial ordenada dá busca binária $\Theta(\log n)$, mas inserção/remoção $\Omega(n)$; lista encadeada insere rápido mas busca em $\Omega(n)$. A ABB combina busca eficiente com estrutura dinâmica.
\begin{defi}
Árvore binária em que, para todo nó $p$: a chave de $p$ é $\ge$ às chaves da subárvore esquerda e $\le$ às da subárvore direita. (O percurso \emph{em ordem} lista as chaves ordenadas.)
\end{defi}
\begin{itemize}
  \item \textbf{Busca:} desce comparando; custo $\Theta(h)$, $h$ = altura.
  \item \textbf{Inserção:} busca a posição e insere como \textbf{folha}: $\Theta(h)$.
  \item \textbf{Remoção} de $p$: (i) folha: remove; (ii) um filho: o filho ocupa o lugar de $p$; (iii) dois filhos: o \textbf{maior da subárvore esquerda} (predecessor: desce à esquerda uma vez e depois sempre à direita) assume o lugar de $p$ (alternativa simétrica: menor da direita). $\Theta(h)$.
  \item $\floor{\log_2 n}+1\le h\le n$ (em níveis). Inserções em ordem crescente/decrescente geram uma ``lista'' de altura $n$ $\Rightarrow$ tudo vira $\Theta(n)$.
\end{itemize}
Árvores completas minimizam a altura, mas mantê-las completas após inserções custa $\Omega(n)$ (reorganização global). Solução: exigir só \textbf{balanceamento}: altura $O(\log n)$, e o mesmo para toda subárvore ($m$ nós $\Rightarrow$ altura $O(\log m)$), com rebalanceamento em $O(\log n)$.

\subsection{Árvores AVL (Adel'son-Vel'skii e Landis, 1962)}
\begin{defi}
Uma árvore binária é \textbf{AVL} se, para todo nó $v$, as alturas das subárvores esquerda e direita diferem de no máximo 1 ($v$ é \emph{regulado}; senão, \emph{desregulado}). Define-se
\[
\bal(v)=h_D(v)-h_E(v)\qquad(\text{direita menos esquerda; } \bal=-1 \text{ ``pesa à esquerda''}),
\]
e $v$ é regulado sse $\bal(v)\in\{-1,0,1\}$. Folhas têm $\bal=0$. Toda árvore completa é AVL (a recíproca é falsa).
\end{defi}

\begin{teo}[AVL é balanceada]\label{teo:avlaltura}
Toda árvore AVL com $n$ nós tem altura $O(\log n)$ (precisamente $h<1{,}44\log_2(n+2)$).
\end{teo}
\begin{proof}
Seja $N(h)$ o número mínimo de nós de uma AVL de altura $h$ (em níveis). $N(1)=1$, $N(2)=2$ e, para $h\ge3$, uma AVL mínima de altura $h$ tem raiz, uma subárvore de altura $h-1$ e a outra de altura $h-2$ (a menor possível), ambas mínimas:
\[
N(h)=1+N(h-1)+N(h-2).
\]
(Essas árvores mínimas são as \textbf{árvores de Fibonacci}; $N(h)=F_{h+2}-1$.) Como $N$ é crescente, $N(h)>2N(h-2)$, donde $N(h)\ge2^{\floor{h/2}}$, isto é, $h\le2\log_2N(h)+1$. Como qualquer AVL de altura $h$ com $n$ nós tem $n\ge N(h)$, $h=O(\log n)$.
\end{proof}
Consequência: busca em $O(\log n)$. Falta mostrar que inserção e remoção, com o rebalanceamento, também são $O(\log n)$.

% ------------------------------------------------------------------
\subsection{Inserção: atualização dos balanços}
Insere-se como numa ABB (a nova chave vira folha; se já existe, para). Depois percorre-se o caminho de volta, do pai $w$ do novo nó até a raiz, atualizando $\bal$. Suponha que a inserção foi na subárvore \textbf{esquerda} de $p$ (o caso direito é simétrico, trocando os sinais):
\begin{center}
\small
\begin{tabular}{@{}cccl@{}}
\toprule
$\bal(p)$ antes & depois & altura de $T(p)$ & ação\\
\midrule
$+1$ & $0$ & não mudou & \textbf{para} (ancestrais não mudam)\\
$0$ & $-1$ & aumentou 1 & continua com o pai de $p$ (se $p$ é a raiz, para)\\
$-1$ & $-2$ & --- & $p$ desregulado: \textbf{rotação}; depois, \textbf{para}\\
\bottomrule
\end{tabular}
\end{center}
Uma inserção aumenta alturas em no máximo 1, logo o desregulado tem $|\bal|=2$. O $p$ a rotacionar é o \textbf{primeiro nó desregulado encontrado ao subir da folha para a raiz} (o mais profundo). Como a rotação devolve à subárvore a altura que ela tinha \emph{antes} da inserção, \textbf{uma inserção faz no máximo uma rotação} (simples ou dupla).

\subsection{As rotações}
Seja $p$ o nó desregulado mais profundo.
\begin{center}
\small
\begin{tabular}{@{}lllll@{}}
\toprule
Caso & $\bal(p)$ & filho do lado pesado & sinais & rotação\\
\midrule
1.1 (``LL'') & $-2$ & $u=$ esq., $\bal(u)=-1$ & iguais & \textbf{simples à direita} em $p$\\
1.2 (``LR'') & $-2$ & $u=$ esq., $\bal(u)=+1$ & opostos & \textbf{dupla à direita}: esq.\ em $u$, dir.\ em $p$\\
2.1 (``RR'') & $+2$ & $z=$ dir., $\bal(z)=+1$ & iguais & \textbf{simples à esquerda} em $p$\\
2.2 (``RL'') & $+2$ & $z=$ dir., $\bal(z)=-1$ & opostos & \textbf{dupla à esquerda}: dir.\ em $z$, esq.\ em $p$\\
\bottomrule
\end{tabular}
\end{center}
Regra mnemônica (monitoria): \textbf{FB positivo $\Rightarrow$ rotações à esquerda; negativo $\Rightarrow$ à direita. Mesmo sinal: simples; sinais opostos: dupla.} (Filho com $\bal=0$ só ocorre na \emph{remoção}, e aí usa-se rotação simples.)

\begin{center}
\begin{tabular}{ccc}
\begin{forest} bst
[$p$ [$u$ [$T_1$,tri] [$T_2$,tri]] [$T_3$,tri]]
\end{forest}
& \raisebox{3em}{$\xrightarrow{\text{rot.\ direita em }p}$} &
\begin{forest} bst
[$u$ [$T_1$,tri] [$p$ [$T_2$,tri] [$T_3$,tri]]]
\end{forest}\\
\multicolumn{3}{c}{\small Caso 1.1: $h(T_1)=h+1$, $h(T_2)=h(T_3)=h$. Depois: $\bal(u)=\bal(p)=0$; altura volta a $h+2$.}\\[6pt]
\begin{forest} bst
[$p$ [$u$ [$T_1$,tri] [$v$ [$T_2$,tri] [$T_3$,tri]]] [$T_4$,tri]]
\end{forest}
& \raisebox{3em}{$\xrightarrow{\text{dupla à direita}}$} &
\begin{forest} bst
[$v$ [$u$ [$T_1$,tri] [$T_2$,tri]] [$p$ [$T_3$,tri] [$T_4$,tri]]]
\end{forest}\\
\multicolumn{3}{c}{\small Caso 1.2: $v$ sobe para a raiz. Depois $\bal(v)=0$; se $\bal(v)$ era $-1$: $\bal(u)=0,\ \bal(p)=+1$;}\\
\multicolumn{3}{c}{\small se era $+1$: $\bal(u)=-1,\ \bal(p)=0$; se era $0$ ($v$ é o novo nó): ambos $0$.}
\end{tabular}
\end{center}
Os casos 2.1 e 2.2 são espelhados (rotação simples à esquerda: $z$ sobe, $p$ vira filho esquerdo de $z$ e o antigo filho esquerdo de $z$ vira filho direito de $p$). \textbf{Rotações preservam a ordem da ABB} (o percurso em ordem $T_1,u,T_2,p,T_3$ é o mesmo) e custam $O(1)$ (só ponteiros).

\subsection{Pseudocódigo (slides do professor)}
Cada nó tem campos \texttt{esq}, \texttt{dir}, \texttt{chave}, \texttt{bal}; $\nil$ é o ponteiro nulo; $pt\ptr.campo$ acessa o campo do nó apontado. Chamada externa: \textsc{Ins-AVL}$(x,ptraiz)$. O booleano $h$ devolvido indica se \textbf{a altura da subárvore aumentou} (só então o pai precisa atualizar seu balanço).

\begin{pseudo}{Inserção}
\begin{algorithmic}
\Procedure{Iniciar-no}{$pt, x$}
  \State alocar$(pt)$;\ $pt\ptr.esq:=\nil$;\ $pt\ptr.dir:=\nil$;\ $pt\ptr.chave:=x$;\ $pt\ptr.bal:=0$
\EndProcedure
\Function{Ins-AVL}{$x$, \textbf{ref} $pt$} \Comment{retorna true se a altura de $T(pt)$ aumentou}
  \If{$pt=\nil$} \State \Call{Iniciar-no}{$pt,x$};\ \Return true \EndIf
  \If{$x=pt\ptr.chave$} \State \Return false \Comment{chave repetida} \EndIf
  \If{$x<pt\ptr.chave$}
    \State $h:=$ \Call{Ins-AVL}{$x, pt\ptr.esq$}
    \If{$h$} \State $pt,h:=$ \Call{AtualizaBalanço}{$pt$, ``esquerda''} \EndIf
  \Else
    \State $h:=$ \Call{Ins-AVL}{$x, pt\ptr.dir$}
    \If{$h$} \State $pt,h:=$ \Call{AtualizaBalanço}{$pt$, ``direita''} \EndIf
  \EndIf
  \State \Return $h$
\EndFunction
\Function{AtualizaBalanço}{$pt, a$}
  \State $h:=$ false
  \If{$a=$ ``esquerda''}
    \If{$pt\ptr.bal=1$} \State $pt\ptr.bal:=0$
    \ElsIf{$pt\ptr.bal=0$} \State $pt\ptr.bal:=-1$;\ $h:=$ true \Comment{altura aumentou: o pai verifica}
    \Else \State $pt:=$ \Call{RotaçõesDireita}{$pt$} \Comment{era $-1$: ficaria $-2$}
    \EndIf
  \Else
    \If{$pt\ptr.bal=-1$} \State $pt\ptr.bal:=0$
    \ElsIf{$pt\ptr.bal=0$} \State $pt\ptr.bal:=1$;\ $h:=$ true
    \Else \State $pt:=$ \Call{RotaçõesEsquerda}{$pt$}
    \EndIf
  \EndIf
  \State \Return $(pt,h)$ \Comment{se houve rotação, $pt$ aponta para a nova raiz da subárvore}
\EndFunction
\end{algorithmic}
\end{pseudo}

\begin{pseudo}{Rotações (slides)}
\begin{algorithmic}
\Function{RotaçõesDireita}{$pt$}
  \State $ptu:=pt\ptr.esq$
  \If{$ptu\ptr.bal=-1$} \Comment{simples à direita}
    \State $pt\ptr.esq:=ptu\ptr.dir$;\ $ptu\ptr.dir:=pt$;\ $pt\ptr.bal:=0$;\ $ptu\ptr.bal:=0$;\ $pt:=ptu$
  \Else \Comment{dupla à direita}
    \State $ptv:=ptu\ptr.dir$
    \State $ptu\ptr.dir:=ptv\ptr.esq$;\ $ptv\ptr.esq:=ptu$;\ $pt\ptr.esq:=ptv\ptr.dir$;\ $ptv\ptr.dir:=pt$
    \If{$ptv\ptr.bal=-1$} $pt\ptr.bal:=1$ \Else\ $pt\ptr.bal:=0$ \EndIf
    \If{$ptv\ptr.bal=1$} $ptu\ptr.bal:=-1$ \Else\ $ptu\ptr.bal:=0$ \EndIf
    \State $ptv\ptr.bal:=0$;\ $pt:=ptv$
  \EndIf
  \State \Return $pt$
\EndFunction
\Function{RotaçõesEsquerda}{$pt$}
  \State $ptz:=pt\ptr.dir$
  \If{$ptz\ptr.bal=1$} \Comment{simples à esquerda}
    \State $pt\ptr.dir:=ptz\ptr.esq$;\ $ptz\ptr.esq:=pt$;\ $pt\ptr.bal:=0$;\ $ptz\ptr.bal:=0$;\ $pt:=ptz$
  \Else \Comment{dupla à esquerda}
    \State $pty:=ptz\ptr.esq$
    \State $ptz\ptr.esq:=pty\ptr.dir$;\ $pty\ptr.dir:=ptz$;\ $pt\ptr.dir:=pty\ptr.esq$;\ $pty\ptr.esq:=pt$
    \If{$pty\ptr.bal=1$} $pt\ptr.bal:=-1$ \Else\ $pt\ptr.bal:=0$ \EndIf
    \If{$pty\ptr.bal=-1$} $ptz\ptr.bal:=1$ \Else\ $ptz\ptr.bal:=0$ \EndIf
    \State $pty\ptr.bal:=0$;\ $pt:=pty$
  \EndIf
  \State \Return $pt$
\EndFunction
\end{algorithmic}
\end{pseudo}

\begin{cuidado}[Detalhes do pseudocódigo que costumam ser perguntados]
\begin{itemize}
  \item \textbf{Para que serve $h$?} Indica se a altura da subárvore cresceu; só nesse caso o pai precisa atualizar o balanço. Por isso \textsc{AtualizaBalanço} só é chamada se $h=$ true: se a altura não mudou, nenhum ancestral muda de balanço.
  \item \textbf{Por que atualizar $pt$ após a rotação?} A raiz da subárvore mudou ($u$, $v$, $z$ ou $y$ sobe); como $pt$ é passado \emph{por referência}, a atribuição religa a nova raiz ao pai (ou atualiza $ptraiz$).
  \item Nos slides, os três testes de \textsc{AtualizaBalanço} aparecem como ``se'' consecutivos; eles precisam ser \textbf{mutuamente exclusivos} (``senão se''), como acima. Com ``se'' em sequência, um nó com $\bal=1$ viraria $0$ e, logo em seguida, o teste ``$\bal=0$'' dispararia de novo (bug do ``aluno'' da monitoria).
  \item Após a rotação na inserção, $h:=$ false (a altura voltou ao valor antigo).
\end{itemize}
\end{cuidado}

% ------------------------------------------------------------------
\subsection{Remoção em AVL}
Remove-se como numa ABB (se $p$ tem dois filhos, o maior da subárvore esquerda o substitui) e sobe-se a partir do nó \emph{fisicamente removido} atualizando balanços. Agora a pergunta é se a altura \textbf{diminuiu}. Remoção na subárvore \textbf{direita} de $p$ (esquerda é simétrico):
\begin{center}
\small
\begin{tabular}{@{}cccl@{}}
\toprule
$\bal(p)$ antes & depois & altura de $T(p)$ & ação\\
\midrule
$+1$ & $0$ & \textbf{diminuiu} & continua com o pai\\
$0$ & $-1$ & não mudou & \textbf{para}\\
$-1$ & $-2$ & --- & rotação à direita (ver abaixo)\\
\bottomrule
\end{tabular}
\end{center}
Na rotação, olha-se o filho esquerdo $u$:
\begin{itemize}
  \item $\bal(u)=-1$: simples à direita; ambos ficam $0$; a altura \textbf{diminuiu} $\Rightarrow$ continua subindo.
  \item $\bal(u)=0$ (\emph{caso exclusivo da remoção}): simples à direita; fica $\bal(p)=-1$, $\bal(u)=+1$; a altura \textbf{não muda} $\Rightarrow$ para.
  \item $\bal(u)=+1$: dupla à direita; altura \textbf{diminuiu} $\Rightarrow$ continua subindo.
\end{itemize}
Por isso a remoção pode precisar de uma rotação em \textbf{cada} ancestral: $O(\log n)$ rotações, cada uma $O(1)$ $\Rightarrow O(\log n)$ no total.

\begin{cuidado}
Os slides implementam a remoção em dois passos (\textsc{Remover-AVL-Passo1}: remoção de ABB que devolve o nó mais profundo afetado e o ajuste $+1/0/-1$; \textsc{Remover-AVL-Passo2}: sobe atualizando com \textsc{AtualizaBalançoApósRemoção}). Ao reutilizar \textsc{RotaçõesDireita} na remoção, o caso $ptu\ptr.bal=0$ cairia na rotação \emph{dupla} (o teste é só ``$=-1$''), o que está errado: com $\bal(u)=0$ deve-se fazer a \textbf{simples} (teste ``$\le0$'') e ajustar $\bal(p)=-1$, $\bal(u)=+1$, devolvendo ``altura não diminuiu''.
\end{cuidado}

\begin{exemplo}[Remoções pequenas (conferidas)]\leavevmode
\begin{itemize}
  \item $20^{+1}(10,\ 40^{0}(30,50))$, remover $10$: $\bal(20)=+2$, filho $40$ com $\bal=0$ $\Rightarrow$ simples à esquerda em 20: $40^{-1}(20^{+1}(\cdot,30),\ 50)$; altura inalterada.
  \item $20^{-1}(10^{+1}(\cdot,15),\ 30)$, remover $30$: $\bal(20)=-2$, filho $10$ com $\bal=+1$ $\Rightarrow$ dupla à direita: $15(10,20)$.
\end{itemize}
\end{exemplo}

% ------------------------------------------------------------------
\subsection{Exemplos resolvidos (conferidos por simulação)}
\begin{exemplo}[Teste 1, Q5]\label{ex:teste1q5}
Árvore inicial e inserções $5, 15, 35, 45$ (balanços em vermelho, $\bal=h_D-h_E$):
\begin{center}
\begin{forest} bst
[50,bb=-1 [30,bb=-1 [20,bb=0 [10,bb=0] [25,bb=0]] [40,bb=0]] [70,bb=0 [60,bb=0] [80,bb=0]]]
\end{forest}
\end{center}
\textbf{Inserir 5} (filho esq.\ de 10). Subindo: $\bal(10)=-1$, $\bal(20)=-1$, $\bal(30)=-2$. \textbf{Primeiro desregulado: 30}; filho esq.\ $20$ com $\bal=-1$ (mesmo sinal) $\Rightarrow$ caso 1.1, \textbf{rotação simples à direita em 30}. A altura da subárvore volta a 3, então 50 fica com $-1$.
\begin{center}
\begin{forest} bst
[50,bb=-1 [20,bb=0 [10,bb=-1 [5,novo,bb=0] [,vazio]] [30,bb=0 [25,bb=0] [40,bb=0]]] [70,bb=0 [60,bb=0] [80,bb=0]]]
\end{forest}
\end{center}
\textbf{Inserir 15} (filho dir.\ de 10): $\bal(10)$ vai de $-1$ a $0$ e a altura de $T(10)$ não muda: \textbf{para}. Sem rotação.
\begin{center}
\begin{forest} bst
[50,bb=-1 [20,bb=0 [10,bb=0 [5,bb=0] [15,novo,bb=0]] [30,bb=0 [25,bb=0] [40,bb=0]]] [70,bb=0 [60,bb=0] [80,bb=0]]]
\end{forest}
\end{center}
\textbf{Inserir 35} (filho esq.\ de 40). Subindo: $\bal(40)=-1$, $\bal(30)=+1$, $\bal(20)=+1$, $\bal(50)=-2$. \textbf{Primeiro desregulado: 50}; filho esq.\ $20$ com $\bal=+1$ (sinais opostos) $\Rightarrow$ caso 1.2, \textbf{rotação dupla à direita}: esquerda em 20, depois direita em 50. O nó $v=30$ sobe para a raiz; como $\bal(30)$ era $+1$: $\bal(20)=-1$, $\bal(50)=0$.
\begin{center}
\begin{forest} bst
[30,bb=0 [20,bb=-1 [10,bb=0 [5,bb=0] [15,bb=0]] [25,bb=0]] [50,bb=0 [40,bb=-1 [35,novo,bb=0] [,vazio]] [70,bb=0 [60,bb=0] [80,bb=0]]]]
\end{forest}
\end{center}
\textbf{Inserir 45} (filho dir.\ de 40): $\bal(40)$ vai de $-1$ a $0$, altura de $T(40)$ inalterada: para. Árvore final:
\begin{center}
\begin{forest} bst
[30,bb=0 [20,bb=-1 [10,bb=0 [5,bb=0] [15,bb=0]] [25,bb=0]] [50,bb=0 [40,bb=0 [35,bb=0] [45,novo,bb=0]] [70,bb=0 [60,bb=0] [80,bb=0]]]]
\end{forest}
\end{center}
\textbf{(c)} Numa ABB comum, inserir códigos em ordem desfavorável (crescente, decrescente, ``zigue-zague'') gera altura $\Theta(n)$ e busca/inserção/remoção em $\Theta(n)$; a AVL garante altura $\le1{,}44\log_2 n$ (Teo.~\ref{teo:avlaltura}), logo $O(\log n)$ \emph{independentemente da ordem de chegada}, pagando só $O(1)$ rotações por inserção.
\end{exemplo}

\begin{exemplo}[Monitoria de revisão, Q5: $40,20,10,30,35$]
$40$; $40^{-1}(20)$; inserir $10$: $\bal(40)=-2$, filho $20$ com $-1$ $\Rightarrow$ simples à direita em 40: $20^{0}(10,40)$. Inserir $30$: $20^{+1}(10,\ 40^{-1}(30,\cdot))$. Inserir $35$ (dir.\ de 30): $\bal(30)=+1$, $\bal(40)=-2$ $\Rightarrow$ primeiro desregulado é \textbf{40}, filho esq.\ $30$ com $+1$ $\Rightarrow$ \textbf{dupla à direita} (esq.\ em 30, dir.\ em 40). Final: $20^{+1}(10^{0},\ 35^{0}(30^{0},40^{0}))$. Rotações: na 3ª inserção (simples, direita) e na 5ª (dupla, direita).
\end{exemplo}

\begin{exemplo}[Monitoria de revisão, Q7: $10,20,30,40,50$ --- e um erro no slide]
Inserir $30$: $\bal(10)=+2$, filho $20$ com $+1$ $\Rightarrow$ simples à esquerda em 10: $20(10,30)$. Inserir $40$: $20^{+1}(10,\ 30^{+1}(\cdot,40))$. Inserir $50$: subindo, $\bal(40)=+1$, $\bal(30)=+2$ $\Rightarrow$ o \textbf{primeiro desregulado é 30} (e não 20!): simples à esquerda em 30. \textbf{Final correto:} $20^{+1}(10^{0},\ 40^{0}(30^{0},50^{0}))$.
\begin{center}
\begin{forest} bst
[20,bb=+1 [10,bb=0] [40,bb=0 [30,bb=0] [50,bb=0]]]
\end{forest}
\end{center}
\textbf{Atenção:} o slide da monitoria rotaciona em 20 e chega a uma árvore com raiz 30 --- está errado: a rotação é feita no desregulado \emph{mais profundo}; após rotacionar em 30, o nó 20 fica com $\bal=+1$ e nem chega a desregular.

\textbf{Inserir 49} na árvore correta: vai para a esquerda de 50; $\bal(50)=-1$, $\bal(40)=+1$, $\bal(20)=+2$ $\Rightarrow$ desregulado 20, filho dir.\ $40$ com $+1$ $\Rightarrow$ \textbf{simples à esquerda em 20}: $40^{0}(20^{0}(10,30),$\allowbreak$\ 50^{-1}(49,\cdot))$. (Partindo da árvore errada do slide, $30(20(10),$\allowbreak$40(\cdot,50))$, o desregulado seria 40 com filho $50$ de $\bal=-1$, e aí sim seria a \textbf{dupla à esquerda} mostrada no slide, com resultado $30(20(10),$\allowbreak$49(40,50))$.)
\end{exemplo}

% ------------------------------------------------------------------
\subsection{Exercícios da lista (Union-Find e AVL)}
\begin{exemplo}[Ex.\ 3 -- remoção com rotação em \emph{todos} os ancestrais]
Use uma \textbf{árvore de Fibonacci} (AVL mínima) em que todo nó interno tem $\bal=-1$ (esquerda mais alta): $T_h$ tem raiz, $T_{h-1}$ à esquerda e $T_{h-2}$ à direita. Removendo a \textbf{folha mais à direita} (caminho raiz $\to T_{h-2}\to T_{h-4}\to\cdots$), cada ancestral perde altura do lado já mais baixo, fica com $\bal=-2$ e tem filho esquerdo com $\bal=-1$ $\Rightarrow$ rotação simples que \emph{diminui} a altura, propagando o problema. Exemplo conferido: $T_7$ com 33 nós (chaves $1..33$ em ordem), raiz 21; remover a folha 33 (ancestrais $32,29,21$) causa rotações em \textbf{32, 29 e 21}. Em geral, $h$ ímpar dá $(h-1)/2$ rotações $=\Theta(\log n)$.
\end{exemplo}
\begin{exemplo}[Ex.\ 4 -- fusão de duas AVL (alto nível)]
\emph{Caso ``junção'' (todas as chaves de $T_1$ $<$ todas as de $T_2$):} retire o mínimo $x$ de $T_2$ (ou o máximo de $T_1$) em $O(\log n)$. Suponha $h(T_1)\ge h(T_2)$: desça pela espinha direita de $T_1$ até um nó $v$ com $h(v)\in\{h(T_2),h(T_2)+1\}$; crie o nó $x$ com filhos $T(v)$ e $T_2$ e coloque-o no lugar de $v$; suba rebalanceando como numa inserção (no máximo uma rotação ajusta tudo). Custo $O(|h_1-h_2|+1)=O(\log n)$.
\emph{Caso geral (chaves intercaladas):} percursos em ordem de ambas ($O(n+m)$), intercalação das duas listas ordenadas ($O(n+m)$) e construção de uma árvore completa a partir da lista ordenada (sempre a mediana como raiz, $O(n+m)$) --- total linear. Alternativa: inserir cada nó da menor na maior, $O(m\log(n+m))$.
\end{exemplo}
\begin{exemplo}[Ex.\ 5 -- AVL como estrutura de conjuntos disjuntos?]
Possível: cada conjunto é uma AVL e cada elemento guarda ponteiro para o pai. \Find$(x)$ sobe até a raiz da sua AVL (que identifica o conjunto): $O(\log n)$. \Union: inserir os elementos da menor AVL na maior, $O(k\log n)$ ($k$ = tamanho da menor); com ``menor na maior'', cada elemento é reinserido $\le\log n$ vezes, total $O(n\log^2n)$. Ganha-se ordem (percurso ordenado, busca por chave), mas é pior que o Union-Find ($\approx O(1)$ amortizado) para \Find/\Union.
\end{exemplo}

\begin{prova}[Heap $\times$ AVL $\times$ ABB (revisão, Q6 e Q14)]
\begin{itemize}
  \item \textbf{Fila de prioridade pura} (sempre atender o de maior prioridade): \textbf{heap} --- máximo em $O(1)$, inserir/remover $O(\log n)$, vetor simples. A AVL também faria tudo em $O(\log n)$, mas não dá o máximo em $O(1)$ sem ponteiro extra e é mais pesada.
  \item \textbf{Busca por chave arbitrária, remoção de elemento específico, percurso em ordem, sucessor/predecessor, intervalos}: \textbf{AVL} (o heap não tem ordem entre irmãos: buscar uma chave é $O(n)$).
  \item \textbf{ABB simples $\times$ AVL}: com inserções/remoções frequentes e ordem de chegada imprevisível, \textbf{AVL} (ABB pode degenerar para $\Theta(n)$).
\end{itemize}
\end{prova}

% ------------------------------------------------------------------
\subsection{Árvores rubro-negras (pesquisa)}\label{sec:rn}
\begin{complemento}[ -- pesquisa proposta em 02/09, entrega 19/09 (Cormen, cap.~13)]
\textbf{Definição:} ABB em que cada nó é \emph{vermelho} ou \emph{preto} (1 bit extra) e: (1) a raiz é preta; (2) as folhas nulas ($\nil$) são pretas; (3) um nó vermelho tem os dois filhos pretos (não há dois vermelhos seguidos); (4) todo caminho de um nó até as folhas nulas descendentes tem o mesmo número de nós pretos (\emph{altura negra} $bh$).

\textbf{Por que é balanceada:} por (4), a subárvore de um nó $x$ tem $\ge2^{bh(x)}-1$ nós internos (indução); por (3), pelo menos metade dos nós de qualquer caminho raiz--folha é preta, então $bh(\text{raiz})\ge h/2$. Logo $n\ge2^{h/2}-1$, isto é, $h\le2\log_2(n+1)$.

\textbf{Inserção:} insere como ABB e pinta o nó $z$ de \emph{vermelho} (não altera alturas negras); se o pai é vermelho, viola (3). Correção olhando o \emph{tio} $y$: (caso 1) tio vermelho $\Rightarrow$ recolore pai e tio de preto e avô de vermelho, e continua do avô; (caso 2) tio preto e $z$ em ``zigue-zague'' $\Rightarrow$ rotação no pai para virar caso 3; (caso 3) tio preto e $z$ ``em linha'' $\Rightarrow$ rotação no avô e troca de cores pai/avô. No fim, raiz preta. No máximo \textbf{2 rotações} por inserção; recolorações $O(\log n)$.

\textbf{Remoção:} remoção de ABB; se o nó retirado era preto, surge um ``preto duplo'' que é resolvido olhando o irmão (4 casos), com no máximo \textbf{3 rotações}.

\textbf{AVL $\times$ rubro-negra:} AVL é mais rigidamente balanceada ($h\le1{,}44\log n$ vs.\ $2\log n$): buscas um pouco mais rápidas. Rubro-negra faz menos rotações em atualizações (remoção $O(1)$ rotações vs.\ $O(\log n)$ na AVL): preferida em bibliotecas (\texttt{std::map}, \texttt{TreeMap}, kernel Linux).
\end{complemento}

\begin{chave}
AVL: $\bal=h_D-h_E\in\{-1,0,1\}$; altura $\le1{,}44\log n$ (árvores de Fibonacci). Inserção: sobe atualizando; $+1\to0$ para; $0\to\pm1$ continua; $\pm1\to\pm2$ rotaciona no \textbf{mais profundo} e para (1 rotação). Sinais iguais: simples; opostos: dupla. Remoção: pode rotacionar em todos os ancestrais; filho com $\bal=0$ $\Rightarrow$ simples e para.
\end{chave}

% ================================================================
\section{Ordenação}\label{sec:ordenacao}
% ================================================================
\textbf{Problema:} dado $S$, obter uma permutação $s_1,\dots,s_n$ com $s_i\le s_{i+1}$.
\begin{itemize}
  \item \textbf{Estável:} elementos de chaves iguais aparecem na saída na mesma ordem relativa da entrada (essencial para o radix sort e para ordenar por várias chaves).
  \item \textbf{In-place:} memória adicional independente de $n$ (ou $O(\log n)$ de pilha). Quicksort e heapsort são in-place; mergesort e counting sort não.
\end{itemize}

\subsection{Ordenação por comparação}\label{sec:sortcomp}
\subsubsection*{Mergesort (intercalação)}
\begin{pseudo}{MergeSort e Intercalação}
\begin{algorithmic}
\Procedure{MergeSort}{$V, ini, fim$}
  \If{$ini<fim$}
    \State $meio:=\floor{(ini+fim)/2}$
    \State \Call{MergeSort}{$V,ini,meio$};\ \Call{MergeSort}{$V,meio+1,fim$}
    \State \Call{Intercala}{$V,ini,meio,fim$}
  \EndIf
\EndProcedure
\Procedure{Intercala}{$V,ini,meio,fim$}
  \State $i:=ini$;\ $j:=meio+1$;\ $k:=1$ \Comment{$W$ = vetor auxiliar}
  \While{$i\le meio$ \textbf{e} $j\le fim$}
    \If{$V[i]\le V[j]$} \State $W[k]:=V[i]$;\ $i:=i+1$ \Comment{``$\le$'' mantém a estabilidade}
    \Else \State $W[k]:=V[j]$;\ $j:=j+1$
    \EndIf
    \State $k:=k+1$
  \EndWhile
  \State copie o que sobrou de $V[i..meio]$ ou de $V[j..fim]$ para o fim de $W$
  \State copie $W[1..fim-ini+1]$ de volta para $V[ini..fim]$
\EndProcedure
\end{algorithmic}
\end{pseudo}
$T(n)=2T(n/2)+\Theta(n)\Rightarrow\Theta(n\log n)$ em todos os casos. \textbf{Estável} (com ``$\le$''), \textbf{não in-place} ($\Theta(n)$ extra). Versão iterativa (\emph{bottom-up}): intercale blocos de tamanho $1,2,4,\dots$ ao longo do vetor.

\subsubsection*{Quicksort}
Se $n\le1$, pronto; escolha um \textbf{pivô} $x$; particione em $S_1=\{w<x\}$ e $S_2=\{w\ge x\}$; ordene recursivamente e concatene.
\begin{pseudo}{Quicksort com partição de Lomuto}
\begin{algorithmic}
\Procedure{QuickSort}{$V,ini,fim$}
  \If{$ini<fim$}
    \State $q:=$ \Call{Particiona}{$V,ini,fim$}
    \State \Call{QuickSort}{$V,ini,q-1$};\ \Call{QuickSort}{$V,q+1,fim$}
  \EndIf
\EndProcedure
\Function{Particiona}{$V,ini,fim$}
  \State $x:=V[fim]$;\ $i:=ini-1$ \Comment{pivô no fim (troque antes se escolher outro)}
  \For{$j:=ini$ \textbf{até} $fim-1$}
    \If{$V[j]<x$} \State $i:=i+1$;\ $V[i]\Leftrightarrow V[j]$ \EndIf
  \EndFor
  \State $V[i+1]\Leftrightarrow V[fim]$;\ \Return $i+1$
\EndFunction
\end{algorithmic}
\end{pseudo}
\begin{itemize}
  \item \textbf{Pior caso} $\Theta(n^2)$: partições sempre $0$ e $n-1$ (ex.: pivô $=$ primeiro elemento e vetor já ordenado ou invertido): $T(n)=T(n-1)+\Theta(n)$.
  \item \textbf{Melhor/médio} $\Theta(n\log n)$. Argumento do slide: com pivô ``aleatório'', em metade dos casos a parte maior tem tamanho $\le 3n/4$; se isso sempre ocorresse, no nível $t$ os subproblemas teriam tamanho $\le(3/4)^tn$, $O(\log n)$ níveis de custo $O(n)$; como ocorre com probabilidade $1/2$, em média a cada 2 níveis o tamanho cai por $3/4$: ainda $O(n\log n)$ esperado.
  \item \textbf{Escolha do pivô:} primeiro elemento (ruim se já ordenado), aleatório, \textbf{mediana de três} (primeiro, meio, último). Usando a \textbf{mediana exata via BFPRT} (§\ref{sec:selecao}) obtém-se $O(n\log n)$ no \emph{pior} caso, mas a constante é alta --- na prática não compensa (resposta a um exercício dos slides).
  \item \textbf{In-place} (partição por trocas; recursão $O(\log n)$ de pilha se recursar primeiro na parte menor). \textbf{Não estável} (a troca de longa distância embaralha iguais); dá para torná-lo estável com partição estável usando $O(n)$ de memória extra.
\end{itemize}

\subsubsection*{Heapsort} Ver §\ref{sec:heap}: $\Theta(n\log n)$ sempre, in-place, não estável.

\subsubsection*{Algoritmos quadráticos (para referência)}
\textbf{Inserção}: insere $V[i]$ na parte já ordenada $V[1..i-1]$ deslocando os maiores; $\Theta(n^2)$ pior, $\Theta(n)$ se já ordenado, estável, in-place, ótimo para $n$ pequeno ou quase ordenado (usado dentro do bucket sort). \textbf{Bolha}: $\Theta(n^2)$, estável.

\subsection{Limite inferior $\Omega(n\log n)$ para ordenação por comparação}
\begin{teo}
Todo algoritmo de ordenação baseado em comparações faz $\Omega(n\log n)$ comparações no pior caso.
\end{teo}
\begin{proof}
Modela-se o algoritmo (para cada $n$) por uma \textbf{árvore de decisão}: cada nó interno é uma comparação $a_i\le a_j$?, com dois filhos (os dois resultados); cada folha é um estado final (uma permutação que ordena a entrada). A altura $h$ é o número de comparações no pior caso. Como qualquer uma das $n!$ permutações pode ser a resposta, há pelo menos $n!$ folhas; uma árvore binária de altura $h$ tem no máximo $2^h$ folhas. Logo
\[
2^h\ge n!\ge n(n-1)\cdots\bigl(\tfrac n2\bigr)\ge\bigl(\tfrac n2\bigr)^{n/2}
\ \Rightarrow\ h\ge\tfrac n2\log_2\tfrac n2=\tfrac n2\log n-\tfrac n2=\Omega(n\log n).\qedhere
\]
\end{proof}
\begin{cor}
Heapsort e mergesort são \textbf{ótimos} entre os algoritmos por comparação (supondo comparação $O(1)$). Para ordenar em $o(n\log n)$ é preciso \emph{não} se basear em comparações: usar a estrutura das chaves (inteiros pequenos, dígitos, distribuição).
\end{cor}

% ------------------------------------------------------------------
\subsection{Ordenação em tempo linear}\label{sec:sortlin}
\textbf{Resposta ao desafio ``dá para ordenar em tempo linear?'':} depende da natureza dos elementos.

\subsubsection*{Counting sort (inteiros em $[0,k]$)}
\begin{pseudo}{COUNTING-SORT$(A,n,k)$ (slides)}
\begin{algorithmic}
\For{$i:=0$ \textbf{até} $k$} \State $C[i]:=0$ \EndFor
\For{$j:=1$ \textbf{até} $n$} \State $C[A[j]]:=C[A[j]]+1$ \EndFor \Comment{$C[i]=\#\{j: A[j]=i\}$}
\For{$i:=1$ \textbf{até} $k$} \State $C[i]:=C[i]+C[i-1]$ \EndFor \Comment{$C[i]=\#\{j: A[j]\le i\}$}
\For{$j:=n$ \textbf{decrescendo até} $1$} \Comment{de trás para frente $\Rightarrow$ estável}
  \State $B[C[A[j]]]:=A[j]$;\ $C[A[j]]:=C[A[j]]-1$
\EndFor
\State \Return $B$
\end{algorithmic}
\end{pseudo}
\begin{itemize}
  \item Tempo e espaço $\Theta(n+k)$; \textbf{linear se $k=O(n)$}. Se $k=\Theta(2^n)$ fica exponencial: é \textbf{pseudopolinomial}.
  \item \textbf{Estável}: a última ocorrência de cada valor vai para a última posição do seu bloco. \textbf{Não in-place} (vetores $B$ e $C$). Não faz nenhuma comparação entre elementos.
\end{itemize}

\subsubsection*{Radix sort (inteiros com $d$ dígitos na base $b$)}
Ordena \textbf{dígito a dígito, do menos significativo (LSD) para o mais significativo}, usando em cada passada um método \textbf{estável}.
\begin{pseudo}{Radix sort com filas (slides)}
\begin{algorithmic}
\For{$i:=1$ \textbf{até} $d$}
  \For{$j:=1$ \textbf{até} $n$}
    \State $k:=$ $i$-ésimo dígito menos significativo de $L[j]$ na base $b$;\ insere $L[j]$ no fim da fila $F[k]$
  \EndFor
  \State $j:=1$
  \For{$k:=0$ \textbf{até} $b-1$}
    \While{$F[k]$ não vazia} \State remove o início de $F[k]$ e põe em $L[j]$;\ $j:=j+1$ \EndWhile
  \EndFor
\EndFor
\end{algorithmic}
\end{pseudo}
\begin{itemize}
  \item \textbf{Por que precisa ser estável?} Depois de ordenar pelo dígito $i$, elementos com o mesmo dígito $i$ devem manter a ordem obtida pelos dígitos $1..i-1$; é o que garante (por indução) que após a passada $i$ a lista está ordenada pelos $i$ últimos dígitos. As filas preservam a ordem naturalmente (alternativa: counting sort por dígito).
  \item Tempo $\Theta(d(n+b))$; com $b$ e $d$ constantes, $\Theta(n)$. Truque: inteiros até $n^c$ na base $b=n$ têm $d=c$ dígitos $\Rightarrow\Theta(cn)$.
  \item \textbf{Exemplo da aula (conferido):} $102,313,690,131,222,670,690,130,079,102$.\\
  Unidades: $690\ 670\ 690\ 130\ 131\ 102\ 222\ 102\ 313\ 079$;\\
  dezenas: $102\ 102\ 313\ 222\ 130\ 131\ 670\ 079\ 690\ 690$;\\
  centenas: $079\ 102\ 102\ 130\ 131\ 222\ 313\ 670\ 690\ 690$.
  \item Chaves de tamanhos variáveis $d_i$ (exercício): agrupe por número de dígitos (counting por tamanho) e aplique o radix de modo que cada número só participe das passadas dos seus próprios dígitos, ou use o MSD/American Flag recursivo: $O(\sum d_i)$. Para \textbf{strings}: ordenar da direita para a esquerda completando as menores com um símbolo ``menor que todos'', ou MSD (primeiro caractere) recursivo.
\end{itemize}

\subsubsection*{Bucket sort (reais uniformes em $[0,1)$)}
Divide $[0,1)$ em $r$ baldes iguais, distribui ($A[i]$ vai para o balde $\floor{r\,A[i]}$), ordena cada balde (ex.: inserção) e concatena. Com entrada uniforme, cada balde tem em média $n/r$ elementos; com $r=\Theta(n)$ o tempo \textbf{esperado} é $\Theta(n)$ (Cormen: $E[n_i^2]=2-1/n$, então $\sum E[O(n_i^2)]=O(n)$). Pior caso: o do algoritmo usado nos baldes (tudo num balde só).

\subsubsection*{American Flag sort}
Radix \textbf{MSD in-place}: conta quantos elementos há de cada dígito (vetor $C$ de tamanho $b$), calcula onde começa cada bloco e \emph{permuta no próprio vetor} cada elemento para o seu bloco (trocas); depois recursão em cada bloco no dígito seguinte. \textbf{Memória auxiliar:} só os contadores, $O(b)$ por nível de recursão, $O(b\cdot d)$ no total (independe de $n$) $\Rightarrow$ \textbf{in-place}. \textbf{Não é estável} (as trocas mudam a ordem de iguais).

% ------------------------------------------------------------------
\subsection{Qual algoritmo de ordenação usar?}\label{sec:escolhasort}
\begin{center}
\small
\begin{tabularx}{\linewidth}{@{}lX@{}}
\toprule
Situação & Escolha\\
\midrule
Inteiros num intervalo pequeno ($k=O(n)$) & \textbf{Counting sort} $\Theta(n+k)$\\
Inteiros/IDs com poucos dígitos ($d$ fixo), $k$ enorme & \textbf{Radix sort} $\Theta(d(n+b))$\\
Reais uniformemente distribuídos em um intervalo & \textbf{Bucket sort} $\Theta(n)$ esperado\\
Chaves arbitrárias, precisa de garantia e/ou estabilidade & \textbf{Mergesort} $\Theta(n\log n)$ (memória $\Theta(n)$)\\
Chaves arbitrárias, memória escassa, garantia & \textbf{Heapsort} $\Theta(n\log n)$ in-place\\
Chaves arbitrárias, desempenho médio na prática & \textbf{Quicksort} (pivô aleatório/mediana de 3)\\
Quase ordenado ou $n$ pequeno & \textbf{Inserção}\\
Dados em lista encadeada / arquivos externos & \textbf{Mergesort}\\
\bottomrule
\end{tabularx}
\end{center}

\begin{prova}[Teste 2, Exercício 3 -- respostas esperadas]
\begin{itemize}
  \item \textbf{A} ($10^7$ notas entre $0$ e $100$): \textbf{Counting sort}. $k=101\ll n$: $\Theta(n+k)=\Theta(n)$, vetor de contagem de 101 posições; comparação daria $\Omega(n\log n)\approx 2{,}3\cdot10^8$ comparações.
  \item \textbf{B} ($10^6$ identificadores de 12 dígitos): \textbf{Radix sort}. Counting é inviável ($k=10^{12}$); radix com $d=12$ e $b=10$ faz $12$ passadas lineares: $\Theta(d(n+b))\approx1{,}2\cdot10^7$ (ou base $10^3$/$10^4$: 4 ou 3 passadas).
  \item \textbf{C} ($10^5$ reais uniformes em $[0,1)$): \textbf{Bucket sort}, $n$ baldes, tempo esperado $\Theta(n)$ (radix/counting não se aplicam a reais).
  \item \textbf{D} ($10^6$ inteiros em $[-10^9,10^9]$): \textbf{Radix sort} após somar $10^9$ (chaves em $[0,2\cdot10^9]$, $\approx31$ bits: 2 passadas na base $2^{16}$, ou 10 na base 10). Counting inviável ($k=2\cdot10^9\gg n$ de memória). Mergesort/quicksort $\Theta(n\log n)$ também seriam aceitáveis, mas o radix é linear.
  \item \textbf{E} ($10^5$ registros, chaves arbitrárias, sem restrição de domínio): só \textbf{comparação}: \textbf{Mergesort} (garantia $\Theta(n\log n)$ e estabilidade, ideal para registros) ou Quicksort (mais rápido na média, in-place, mas $\Theta(n^2)$ no pior caso e não estável).
\end{itemize}
\end{prova}

% ------------------------------------------------------------------
\subsection{Q4: dado um código, qual estrutura e qual ordenação?}\label{sec:escolhageral}
A Q4 dará um código genérico; o roteiro é: \textbf{(1)} liste as operações que o código faz repetidamente; \textbf{(2)} estime quantas vezes cada uma ocorre; \textbf{(3)} escolha a estrutura que torna barata a operação dominante; \textbf{(4)} olhe o tipo das chaves para escolher a ordenação; \textbf{(5)} justifique com complexidades (e memória/estabilidade quando relevante).
\begin{center}
\small
\begin{tabularx}{\linewidth}{@{}XX@{}}
\toprule
O código faz repetidamente\dots & Estrutura\\
\midrule
pega o maior/menor, insere, (às vezes) altera prioridade & \textbf{heap binário}\\
\dots\ e une filas com frequência (previsível) & \textbf{heap binomial}\\
\dots\ com muitíssimos aumentos de prioridade (custo total) & \textbf{heap de Fibonacci}\\
``$x$ e $y$ estão no mesmo grupo?'' / une grupos & \textbf{Union-Find} (rank + compressão)\\
busca por chave, inserção, remoção, em ordem, sucessor, mínimo/máximo, intervalos & \textbf{AVL} (ou rubro-negra)\\
só busca/inserção/remoção por chave exata, sem ordem & \textbf{tabela hash} ($O(1)$ esperado)\\
muitas buscas, quase nenhuma atualização & \textbf{vetor ordenado} + busca binária\\
$k$-ésimo menor / mediana uma vez & \textbf{seleção linear} (BFPRT/quickselect), sem ordenar\\
os $k$ maiores de um fluxo & heap de mínimo de tamanho $k$: $O(n\log k)$\\
\bottomrule
\end{tabularx}
\end{center}
\textbf{Padrões a reconhecer em código:} laço que procura o máximo de uma lista e o remove (\,$\Theta(n)$ por iteração $\Rightarrow$ heap); laço que ordena a lista inteira a cada inserção ($\Rightarrow$ heap ou AVL); laço que percorre todos os elementos para saber se dois estão no mesmo grupo ($\Rightarrow$ Union-Find); ordenar para depois pegar só o $k$-ésimo ($\Rightarrow$ seleção $O(n)$); ordenar inteiros pequenos com \texttt{sort} genérico ($\Rightarrow$ counting/radix).

\begin{chave}
Comparação $\Rightarrow\Omega(n\log n)$ (árvore de decisão, $2^h\ge n!$). Merge: estável, $\Theta(n)$ extra. Quick: in-place, médio $n\log n$, pior $n^2$. Heap: in-place, sempre $n\log n$, não estável. Counting $\Theta(n+k)$ estável; Radix $\Theta(d(n+b))$ precisa de estável por dígito (LSD); Bucket $\Theta(n)$ esperado com uniforme.
\end{chave}

% ================================================================
\section{Seleção: o algoritmo BFPRT (mediana das medianas)}\label{sec:selecao}
% ================================================================
\textbf{Problema da seleção:} dado um conjunto $A$ de $n$ números e $k$, achar o $k$-ésimo menor. Casos: mínimo ($k=1$), máximo ($k=n$), mediana inferior $\floor{(n+1)/2}$, superior $\ceil{(n+1)/2}$.

\paragraph{Possibilidades.} (1) Ordenar: $O(n\log n)$. (2) Retirar o mínimo $k$ vezes: $O(kn)$ ($O(n^2)$ para a mediana). (3) \textbf{Quickselect} (pivoteamento como no quicksort, mas recursão em \emph{um} lado só): $O(n)$ esperado, $O(n^2)$ no pior caso. (4) \textbf{BFPRT} (Blum, Floyd, Pratt, Rivest, Tarjan): $O(n)$ no \textbf{pior caso}.

\begin{pseudo}{Quickselect (exercício do slide: tempo esperado $O(n)$, mais simples)}
\begin{algorithmic}
\Function{QuickSelect}{$S,k$}
  \State escolha um pivô $m$ \textbf{aleatório} em $S$
  \State $S_1:=\{s<m\}$;\ $S_2:=\{s=m\}$;\ $S_3:=\{s>m\}$
  \If{$|S_1|\ge k$} \State \Return \Call{QuickSelect}{$S_1,k$}
  \ElsIf{$|S_1|+|S_2|\ge k$} \State \Return $m$
  \Else \State \Return \Call{QuickSelect}{$S_3,\,k-|S_1|-|S_2|$}
  \EndIf
\EndFunction
\end{algorithmic}
\end{pseudo}
Esperado: com probabilidade $\ge1/2$ o pivô cai no ``meio'' (entre os quartis) e o lado mantido tem $\le 3n/4$. Agrupe as chamadas em \emph{fases}: a fase termina quando o tamanho cai para $\le3n/4$. O número de chamadas numa fase é dominado por uma geométrica de parâmetro $1/2$, logo tem esperança $\le2$, e cada chamada custa $\le cn$. Assim $E[T(n)]\le E[T(3n/4)]+2cn$, e somando a série geométrica, $E[T(n)]\le 2cn\sum_{t\ge0}(3/4)^t=8cn=O(n)$.

\subsection{O algoritmo}
A ideia é garantir um pivô \emph{sempre} bom: a \textbf{mediana das medianas} de grupos de 5.
\begin{pseudo}{Seleção$(S,k)$ -- BFPRT (slides)}
\begin{algorithmic}
\Function{Seleção}{$S,k$}
  \If{$|S|\le49$} \State ordene $S$ (mergesort) e \Return o $k$-ésimo menor \EndIf
  \State divida $S$ em $\ceil{|S|/5}$ grupos de (até) 5 elementos
  \State determine a mediana de cada grupo (ordenar 5 elementos: $O(1)$ cada)
  \State $M:=$ conjunto das medianas;\quad $m:=$ \Call{Seleção}{$M,\ceil{|M|/2}$} \Comment{recursão 1: tamanho $n/5$}
  \State $S_1:=\{s\in S: s<m\}$;\ $S_2:=\{s\in S: s=m\}$;\ $S_3:=\{s\in S: s>m\}$
  \If{$|S_1|\ge k$} \State \Return \Call{Seleção}{$S_1,k$} \Comment{recursão 2: tamanho $\le3n/4$}
  \ElsIf{$|S_1|+|S_2|\ge k$} \State \Return $m$
  \Else \State \Return \Call{Seleção}{$S_3,\ k-|S_1|-|S_2|$}
  \EndIf
\EndFunction
\end{algorithmic}
\end{pseudo}
\begin{cuidado}
Num dos slides o teste aparece como ``se $|S_1|=k$ então o $k$-ésimo está em $S_1$''; o correto é \textbf{$|S_1|\ge k$} (como no pseudocódigo da página seguinte). E, ao ir para $S_3$, o índice muda para $k-|S_1|-|S_2|$.
\end{cuidado}

\subsection{Análise}
\begin{lema}
$|S_1|\le n-3\floor{n/10}$ e $|S_3|\le n-3\floor{n/10}$; para $n\ge50$, isso é $\le 3n/4$.
\end{lema}
\begin{proof}
Das $\floor{n/5}$ medianas, pelo menos metade, $\floor{\floor{n/5}/2}=\floor{n/10}$, são $\ge m$. Cada uma dessas vem de um grupo que tem mais 2 elementos $\ge$ ela (logo $\ge m$). Assim, pelo menos $3\floor{n/10}$ elementos são $\ge m$, ou seja, $|S_2|+|S_3|\ge3\floor{n/10}$ e $|S_1|\le n-3\floor{n/10}$. Simetricamente para $S_3$.
Para $n-3\floor{n/10}\le3n/4$ basta $12\floor{n/10}\ge n$, isto é, $2\floor{n/10}\ge n-10\floor{n/10}$; como o lado direito é $\le9$, basta $\floor{n/10}\ge5$, ou seja, $n\ge50$ (por isso o caso base é $n\le49$).
\end{proof}

\begin{teo}
$T(n)\le 20cn$, logo a BFPRT é $O(n)$ no pior caso.
\end{teo}
\begin{proof}
$T(n)=O(1)$ para $n\le49$ e $T(n)\le T(n/5)+T(3n/4)+cn$ para $n>49$ (a partição e as medianas dos grupos são lineares). Por indução forte (com a base ajustada, p.ex.\ $T(n)\le6n\le20cn$ pelo mergesort para $n<50$):
\[
T(n)\le 20c\,\frac n5+20c\,\frac{3n}{4}+cn=4cn+15cn+cn=20cn.\qedhere
\]
\end{proof}
\paragraph{Por que grupos de 5?} Para linearidade, a soma dos tamanhos dos subproblemas precisa ser uma fração $<1$ de $n$: $\frac n5+\frac{3n}4=\frac{19n}{20}<n$. Com grupos de 3: pivô garante só $\approx n/3$ de cada lado, $T(n)=T(n/3)+T(2n/3)+cn$ e $\frac13+\frac23=1\Rightarrow\Theta(n\log n)$. Grupos de 7, 9, \dots\ também funcionam (constantes diferentes).

\paragraph{Aplicação.} Quicksort com a mediana (BFPRT) como pivô: $T(n)=2T(n/2)+O(n)=O(n\log n)$ no \textbf{pior} caso; na prática é mais lento que o quicksort aleatório por causa das constantes.

\begin{chave}
Seleção em $O(n)$: grupos de 5, mediana das medianas como pivô, $\ge3\floor{n/10}$ de cada lado, $T(n)\le T(n/5)+T(3n/4)+cn\le20cn$. Não é preciso ordenar para achar o $k$-ésimo!
\end{chave}

% ================================================================
\section{Divisão e conquista}\label{sec:dc}
% ================================================================
\textbf{Paradigma:} (1) \emph{Divisão}: quebrar o problema $\Pi$ em subproblemas menores \textbf{do mesmo tipo} $\Pi_1,\dots,\Pi_k$; (2) \emph{Conquista}: resolvê-los recursivamente (casos pequenos diretamente); (3) \emph{Combinação}: juntar as soluções parciais com um procedimento ``simples''. Busca-se \textbf{bom balanceamento} (partes de tamanhos parecidos). Funciona bem quando os subproblemas são \textbf{disjuntos}; quando eles se repetem muito, a técnica certa é programação dinâmica (§\ref{sec:pd}).
Exemplos do curso: mergesort, quicksort, BFPRT, contagem de inversões, Strassen, subvetor máximo.

\begin{prova}[Q5: modelo de resposta para ``escreva o pseudocódigo por divisão e conquista'']
\begin{enumerate}
  \item \textbf{Assinatura}: \textsc{Resolve}$(V,ini,fim)$ trabalhando num intervalo (evita copiar vetores). Diga o que a função \textbf{retorna} (às vezes mais de um valor: soma \emph{e} índices; contagem \emph{e} vetor ordenado).
  \item \textbf{Caso base}: $ini=fim$ (ou $n\le$ constante): resposta direta.
  \item \textbf{Divisão}: $meio:=\floor{(ini+fim)/2}$.
  \item \textbf{Conquista}: duas chamadas recursivas, $[ini,meio]$ e $[meio+1,fim]$.
  \item \textbf{Combinação}: o que ``atravessa'' o meio (a parte criativa); geralmente $O(n)$.
  \item \textbf{Complexidade}: escreva a recorrência e resolva (árvore de recursão ou Teorema Mestre, §\ref{sec:recorrencias}): $2T(n/2)+\Theta(n)=\Theta(n\log n)$.
  \item \textbf{Correção} (uma frase): toda solução está inteira à esquerda, inteira à direita ou cruza o meio; cada caso é tratado.
\end{enumerate}
\end{prova}

% ------------------------------------------------------------------
\subsection{Contagem de inversões (aula de 11/09)}
Dada $a_1,\dots,a_n$, um par $i<j$ é uma \textbf{inversão} se $a_i>a_j$ (métrica de similaridade entre rankings: minha ordem $1,\dots,n$ vs.\ a sua). Força bruta: $\Theta(n^2)$.
\textbf{D\&C:} divide em metades $A$ e $B$; conta recursivamente as inversões dentro de cada uma; \textbf{combina} contando os pares $(a,b)$ com $a\in A$, $b\in B$, $a>b$. Se $A$ e $B$ estão \emph{ordenadas}, isso é fácil durante a intercalação: ao comparar $a_i$ e $b_j$, se $a_i\le b_j$, $a_i$ não é invertido com nenhum elemento restante de $B$; se $a_i>b_j$, então $b_j$ é invertido com \textbf{todos os elementos restantes de $A$}. Por isso a função também devolve o vetor ordenado (é um mergesort ``instrumentado'').
\begin{pseudo}{Contagem de inversões (aula + completado)}
\begin{algorithmic}
\Function{CalculaInversões}{$V,ini,fim$} \Comment{conta e deixa $V[ini..fim]$ ordenado}
  \If{$ini=fim$} \State \Return $0$ \EndIf
  \State $meio:=\floor{(ini+fim)/2}$
  \State $a:=$ \Call{CalculaInversões}{$V,ini,meio$}
  \State $b:=$ \Call{CalculaInversões}{$V,meio+1,fim$}
  \State $c:=$ \Call{MergeCalculandoInversões}{$V,ini,meio,fim$}
  \State \Return $a+b+c$
\EndFunction
\Function{MergeCalculandoInversões}{$V,ini,meio,fim$}
  \State $i:=ini$;\ $j:=meio+1$;\ $k:=1$;\ $cont:=0$
  \While{$i\le meio$ \textbf{e} $j\le fim$}
    \If{$V[i]\le V[j]$} \State $W[k]:=V[i]$;\ $i:=i+1$
    \Else
      \State $W[k]:=V[j]$;\ $j:=j+1$
      \State $cont:=cont+(meio-i+1)$ \Comment{$V[j]$ é menor que todos os restantes da esquerda}
    \EndIf
    \State $k:=k+1$
  \EndWhile
  \State copie os restantes de $V[i..meio]$ e de $V[j..fim]$ para $W$;\ copie $W$ para $V[ini..fim]$
  \State \Return $cont$
\EndFunction
\end{algorithmic}
\end{pseudo}
$T(n)=2T(n/2)+\Theta(n)\Rightarrow\Theta(n\log n)$. \emph{Cuidado:} usar ``$\le$'' (iguais não são inversão) e somar $meio-i+1$ (quantos restam em $A$), não 1. Ex.\ conferido: $1,5,4,8,10,2,6,9,12,11,3,7$ tem \textbf{22} inversões.

% ------------------------------------------------------------------
\subsection{Subvetor contíguo de soma máxima (Teste 2, Exercício 1)}
Dado $A[1..n]$ (inteiros, podem ser negativos), achar $i\le j$ maximizando $A[i]+\dots+A[j]$ (subvetor \textbf{não vazio}).
\textbf{Ideia:} o melhor subvetor de $A[ini..fim]$ está (i) inteiro na metade esquerda, (ii) inteiro na direita, ou (iii) \textbf{cruza o meio}, i.e., é $A[i..meio]\cup A[meio+1..j]$. O caso (iii) se resolve em tempo linear: melhor sufixo da esquerda terminando em $meio$ + melhor prefixo da direita começando em $meio+1$.
\begin{pseudo}{Subvetor de soma máxima por D\&C (retorna soma, início, fim)}
\begin{algorithmic}
\Function{MaxSub}{$A,ini,fim$}
  \If{$ini=fim$} \State \Return $(A[ini],\,ini,\,ini)$ \Comment{não vazio: um elemento}
  \EndIf
  \State $meio:=\floor{(ini+fim)/2}$
  \State $(s_E,i_E,j_E):=$ \Call{MaxSub}{$A,ini,meio$}
  \State $(s_D,i_D,j_D):=$ \Call{MaxSub}{$A,meio+1,fim$}
  \State $(s_C,i_C,j_C):=$ \Call{MaxCruzado}{$A,ini,meio,fim$}
  \State \Return a tripla de maior soma entre $(s_E,i_E,j_E)$, $(s_D,i_D,j_D)$, $(s_C,i_C,j_C)$
\EndFunction
\Function{MaxCruzado}{$A,ini,meio,fim$}
  \State $soma:=0$;\ $melhorE:=-\infty$
  \For{$i:=meio$ \textbf{decrescendo até} $ini$}
    \State $soma:=soma+A[i]$
    \If{$soma>melhorE$} \State $melhorE:=soma$;\ $iMax:=i$ \EndIf
  \EndFor
  \State $soma:=0$;\ $melhorD:=-\infty$
  \For{$j:=meio+1$ \textbf{até} $fim$}
    \State $soma:=soma+A[j]$
    \If{$soma>melhorD$} \State $melhorD:=soma$;\ $jMax:=j$ \EndIf
  \EndFor
  \State \Return $(melhorE+melhorD,\ iMax,\ jMax)$
\EndFunction
\end{algorithmic}
\end{pseudo}
\begin{itemize}
  \item \textbf{Complexidade:} $T(n)=2T(n/2)+\Theta(n)\Rightarrow\Theta(n\log n)$ (Cormen §4.1).
  \item \textbf{Todos negativos:} funciona, pois o caso base devolve um elemento e o cruzado usa pelo menos $A[meio]$ e $A[meio+1]$; o máximo das três opções é o maior elemento. Ex.\ 2 do teste: $[-8,-3,-6,-2,-5,-4]\to -2$.
  \item Ex.\ 1 do teste (conferido): $[-2,1,-3,4,-1,2,1,-5,4]\to$ soma $6$, subvetor $[4,-1,2,1]$. \emph{O enunciado indexa a partir de 0} (posições 3 a 6); com índices a partir de 1, posições 4 a 7.
  \item (Fora de D\&C: o algoritmo de Kadane resolve em $\Theta(n)$ mantendo o melhor sufixo terminando em cada posição. Também existe D\&C linear devolvendo (soma total, melhor prefixo, melhor sufixo, melhor) de cada metade: $T(n)=2T(n/2)+O(1)=\Theta(n)$.)
\end{itemize}

% ------------------------------------------------------------------
\subsection{Multiplicação de matrizes: Strassen (Teste 2, Exercício 2)}
Produto por blocos ingênuo: 8 multiplicações de $n/2\times n/2$, $T(n)=8T(n/2)+\Theta(n^2)=\Theta(n^3)$. Strassen usa \textbf{7}: $T(n)=7T(n/2)+\Theta(n^2)=\Theta(n^{\log_27})\approx\Theta(n^{2{,}807})$ (Teorema Mestre caso 1, pois $n^2=O(n^{\log_27-\varepsilon})$).
\begin{pseudo}{Strassen$(A,B,n)$ -- $n$ potência de 2; Soma/Subtrai são dadas}
\begin{algorithmic}
\Function{Strassen}{$A,B,n$}
  \If{$n=1$} \State \Return $[\,A[1,1]\cdot B[1,1]\,]$ \EndIf
  \State particione $A=\begin{bmatrix}A_{11}&A_{12}\\A_{21}&A_{22}\end{bmatrix}$, $B=\begin{bmatrix}B_{11}&B_{12}\\B_{21}&B_{22}\end{bmatrix}$ em blocos $\frac n2\times\frac n2$
  \State $M_1:=$ \Call{Strassen}{$A_{11}+A_{22},\ B_{11}+B_{22},\ n/2$}
  \State $M_2:=$ \Call{Strassen}{$A_{21}+A_{22},\ B_{11},\ n/2$}
  \State $M_3:=$ \Call{Strassen}{$A_{11},\ B_{12}-B_{22},\ n/2$}
  \State $M_4:=$ \Call{Strassen}{$A_{22},\ B_{21}-B_{11},\ n/2$}
  \State $M_5:=$ \Call{Strassen}{$A_{11}+A_{12},\ B_{22},\ n/2$}
  \State $M_6:=$ \Call{Strassen}{$A_{21}-A_{11},\ B_{11}+B_{12},\ n/2$}
  \State $M_7:=$ \Call{Strassen}{$A_{12}-A_{22},\ B_{21}+B_{22},\ n/2$}
  \State $C_{11}:=M_1+M_4-M_5+M_7$;\quad $C_{12}:=M_3+M_5$
  \State $C_{21}:=M_2+M_4$;\quad $C_{22}:=M_1-M_2+M_3+M_6$
  \State \Return $\begin{bmatrix}C_{11}&C_{12}\\C_{21}&C_{22}\end{bmatrix}$
\EndFunction
\end{algorithmic}
\end{pseudo}
São 18 somas/subtrações de matrizes $\frac n2\times\frac n2$ ($\Theta(n^2)$) por chamada. Na prática usa-se o método clássico abaixo de um certo $n$ (constantes maiores, menos estabilidade numérica).

% ------------------------------------------------------------------
\subsection{Outros exemplos clássicos de D\&C (treino para a Q5)}
\begin{pseudo}{Busca binária, potência rápida e máximo/mínimo}
\begin{algorithmic}
\Function{BuscaBin}{$V,ini,fim,x$} \Comment{$V$ ordenado; $T(n)=T(n/2)+O(1)=O(\log n)$}
  \If{$ini>fim$} \State \Return $-1$ \EndIf
  \State $meio:=\floor{(ini+fim)/2}$
  \If{$V[meio]=x$} \State \Return $meio$
  \ElsIf{$x<V[meio]$} \State \Return \Call{BuscaBin}{$V,ini,meio-1,x$}
  \Else \State \Return \Call{BuscaBin}{$V,meio+1,fim,x$}
  \EndIf
\EndFunction
\Function{Pot}{$a,n$} \Comment{$a^n$ com $O(\log n)$ multiplicações}
  \If{$n=0$} \State \Return $1$ \EndIf
  \State $t:=$ \Call{Pot}{$a,\floor{n/2}$}
  \If{$n$ par} \State \Return $t\cdot t$ \Else \State \Return $t\cdot t\cdot a$ \EndIf
\EndFunction
\Function{MaxMin}{$V,ini,fim$} \Comment{$\approx 3n/2$ comparações (contra $2n-2$ do ingênuo)}
  \If{$ini=fim$} \State \Return $(V[ini],V[ini])$ \EndIf
  \If{$fim=ini+1$} \State \Return $(\max,\min)$ de $V[ini],V[fim]$ com 1 comparação \EndIf
  \State $meio:=\floor{(ini+fim)/2}$
  \State $(M_1,m_1):=$ \Call{MaxMin}{$V,ini,meio$};\ $(M_2,m_2):=$ \Call{MaxMin}{$V,meio+1,fim$}
  \State \Return $(\max\{M_1,M_2\},\ \min\{m_1,m_2\})$
\EndFunction
\end{algorithmic}
\end{pseudo}

\begin{pseudo}{Elemento majoritário (aparece $>n/2$ vezes), $T(n)=2T(n/2)+\Theta(n)=\Theta(n\log n)$}
\begin{algorithmic}
\Function{Maj}{$V,ini,fim$} \Comment{retorna o majoritário de $V[ini..fim]$ ou \textsc{nenhum}}
  \If{$ini=fim$} \State \Return $V[ini]$ \EndIf
  \State $meio:=\floor{(ini+fim)/2}$
  \State $a:=$ \Call{Maj}{$V,ini,meio$};\ $b:=$ \Call{Maj}{$V,meio+1,fim$}
  \For{$c\in\{a,b\}$, $c\ne$ \textsc{nenhum}}
    \If{$c$ ocorre mais de $(fim-ini+1)/2$ vezes em $V[ini..fim]$} \State \Return $c$ \EndIf \Comment{contagem $O(n)$}
  \EndFor
  \State \Return \textsc{nenhum}
\EndFunction
\end{algorithmic}
\end{pseudo}
\emph{Correção:} se $x$ é majoritário no todo, é majoritário em pelo menos uma das metades (senão ocorreria $\le n/2$ vezes).

\paragraph{Mais ideias de D\&C.} \emph{Karatsuba} (inteiros de $n$ dígitos com 3 produtos de $n/2$: $\Theta(n^{\log_23})\approx n^{1{,}585}$); \emph{par de pontos mais próximos} no plano ($\Theta(n\log n)$: divide pela mediana de $x$, combina olhando a faixa central com no máximo 7 vizinhos por ponto); \emph{pico em vetor unimodal} e \emph{raiz de função monótona} (busca binária, $O(\log n)$); \emph{contar ocorrências de $x$ em vetor ordenado} (duas buscas binárias).

\begin{cuidado}[Erros frequentes em pseudocódigo de D\&C]
\begin{itemize}
  \item Esquecer o caso base ou fazê-lo errado (ex.: subvetor \emph{vazio} permitido quando o enunciado pede não vazio).
  \item Dividir em $[ini,meio-1]$ e $[meio,fim]$ com $meio=\floor{(ini+fim)/2}$: com 2 elementos gera recursão infinita. Use $[ini,meio]$ e $[meio+1,fim]$.
  \item Combinação que olha só um elemento de cada lado (o subvetor cruzado precisa do \emph{melhor sufixo} e do \emph{melhor prefixo}).
  \item Não devolver a informação extra necessária (índices; vetor ordenado nas inversões).
  \item Não dizer a complexidade (o enunciado sempre pede!).
\end{itemize}
\end{cuidado}

\begin{chave}
D\&C = dividir (meio), conquistar (2 chamadas), combinar (o que cruza o meio). $2T(n/2)+\Theta(n)=\Theta(n\log n)$; $T(n/2)+O(1)=\Theta(\log n)$; $7T(n/2)+\Theta(n^2)=\Theta(n^{2{,}807})$; $T(n/5)+T(3n/4)+\Theta(n)=\Theta(n)$.
\end{chave}

% ================================================================
\section{Complementos}
% ================================================================

\subsection{Tabelas hash (pesquisa)}\label{sec:hash}
\begin{complemento}[ -- pesquisa proposta em 02/09, entrega 19/09 (Cormen, cap.~11; Szwarcfiter, cap.~10)]
\textbf{Motivação:} implementar um \emph{dicionário} (inserir, buscar, remover por chave) com acesso ``direto''. Endereçamento direto (vetor indexado pela chave) é $O(1)$ mas exige memória do tamanho do universo de chaves (infinito ou enorme). Uma \textbf{função hash} $h:U\to\{0,\dots,m-1\}$ transforma a chave num endereço de uma tabela de $m$ posições.

\textbf{Funções clássicas:} \emph{divisão} $h(k)=k\bmod m$ ($m$ primo longe de potências de 2); \emph{multiplicação} $h(k)=\floor{m\,(kA\bmod1)}$ ($A\approx(\sqrt5-1)/2$); \emph{dobra}/\emph{meio do quadrado}; para strings, polinomial $h(s)=\sum s_i\,b^{i}\bmod m$; \emph{hashing universal} (escolher $h$ aleatoriamente de uma família, p.ex.\ $((ak+b)\bmod p)\bmod m$), que garante bom caso esperado para qualquer entrada.

\textbf{Colisões} ($h(k_1)=h(k_2)$, $k_1\ne k_2$) são inevitáveis (princípio da casa dos pombos: $|U|>m$). Tratamento:
\begin{itemize}
  \item \textbf{Encadeamento exterior:} cada posição aponta para uma lista das chaves com aquele endereço.
  \item \textbf{Encadeamento interior:} as listas usam posições da própria tabela (uma zona de colisões, ou ponteiros entre posições livres).
  \item \textbf{Endereçamento aberto:} tudo fica na tabela; em colisão, sonda-se uma sequência de posições: \emph{linear} $h(k)+i$ (agrupamento primário), \emph{quadrática} $h(k)+c_1i+c_2i^2$, \emph{duplo hashing} $h_1(k)+i\,h_2(k)$. Remoção exige marca ``removido''.
\end{itemize}
\textbf{Complexidade:} com fator de carga $\alpha=n/m$ e hash uniforme, busca em $\Theta(1+\alpha)$ esperado com encadeamento, e $\le\frac1{1-\alpha}$ sondagens esperadas (busca sem sucesso) no endereçamento aberto; mantendo $\alpha=O(1)$ (redimensionando a tabela, custo amortizado $O(1)$), as operações são $O(1)$ esperado. \textbf{Pior caso: $\Theta(n)$} (todas as chaves no mesmo endereço). Hash não mantém ordem (sem mínimo, sucessor, intervalos --- aí usa-se AVL/rubro-negra).
\end{complemento}

\subsection{Programação dinâmica}\label{sec:pd}
\begin{complemento}[ -- slides de PD (subsequência comum máxima); não previsto no formato da AV1]
Na decomposição recursiva, \textbf{subproblemas idênticos} podem aparecer muitas vezes, gerando tempo exponencial mesmo com poucos subproblemas distintos (ex.: Fibonacci recursivo). \textbf{PD}: guardar a solução de cada subproblema numa tabela e resolvê-los de baixo para cima (\emph{bottom-up}) ou com memorização. D\&C é bom quando os subproblemas são disjuntos; PD quando eles se sobrepõem.

\textbf{Subsequência comum máxima (LCS).} $X=x_1..x_n$, $Y=y_1..y_m$; $C(i,j)=$ tamanho da LCS de $X_i$ e $Y_j$:
\[
C(i,j)=\begin{cases}0 & i=0\text{ ou }j=0\\ 1+C(i-1,j-1) & x_i=y_j\\ \max\{C(i-1,j),C(i,j-1)\} & x_i\ne y_j\end{cases}
\]
Preenche-se a tabela linha a linha: $O(nm)$ tempo e espaço; resposta $C(n,m)$. Ex.: $X=5,2,1,3,9,10,14$, $Y=2,3,1,10$: LCS de tamanho 3 ($2,3,10$ ou $2,1,10$). Para \emph{recuperar} a subsequência, volte de $(n,m)$: se $x_i=y_j$, $x_i$ faz parte e vá a $(i-1,j-1)$; senão vá para a vizinha de maior valor.
\end{complemento}

% ================================================================
\section{Banco de questões com respostas}\label{sec:banco}
% ================================================================

\subsection{Testes aplicados}\label{sec:testes}
\paragraph{Teste 1 (heaps, Union-Find, AVL).}
\begin{enumerate}
  \item Aeroporto: filas de prioridade por setor, fusões muito frequentes, alterações raras, tempo \emph{previsível por operação}. Qual heap? $\Rightarrow$ \textbf{binomial} (§\ref{sec:escolhaheap}).
  \item Central médica: $2\cdot10^6$ chamadas de 16 bytes, limite de 40 MB, ponteiros de 8 bytes, sem fusão, poucas alterações. $\Rightarrow$ \textbf{heap binário} (§\ref{sec:escolhaheap}, conta de memória).
  \item Supercomputador: $5\cdot10^6$ inserções/remoções e $8\cdot10^7$ aumentos de prioridade, custo total importa. $\Rightarrow$ \textbf{Fibonacci} (§\ref{sec:escolhaheap}).
  \item Plataforma com $10^7$ usuários: $10^5$ uniões e $10^8$ consultas; árvore vs.\ lista; estratégia de união. $\Rightarrow$ Union-Find, lista com ponteiro ao representante (\Find{} $O(1)$), união ``menor na maior'' (§\ref{sec:ufescolha}).
  \item AVL: árvore $50(30(20(10,25),40),70(60,80))$, inserir $5,15,35,45$, identificar desregulados e rotações; ABB vs.\ AVL $\Rightarrow$ Exemplo \ref{ex:teste1q5}.
\end{enumerate}
\paragraph{Teste 2 (ordenação e D\&C, 2h, 3 questões).}
\begin{enumerate}
  \item Subvetor contíguo de soma máxima por D\&C, com posições e complexidade $\Rightarrow$ §\ref{sec:dc}.
  \item Strassen: algoritmo recursivo com as 7 multiplicações $\Rightarrow$ §\ref{sec:dc}.
  \item Escolher Counting/Radix/Bucket/Merge/Quick para 5 cenários $\Rightarrow$ §\ref{sec:escolhasort}.
\end{enumerate}

% ------------------------------------------------------------------
\subsection{Lista -- heaps binários, binomiais e de Fibonacci}
\emph{Obs.:} a numeração segue a da lista, que \textbf{pula o exercício 9} (vai do 8 ao 10).
\begin{enumerate}
  \item \emph{Verificar se são heaps.} (a) $33\,32\,28\,31\,26\,29\,25\,30\,27$: \textbf{não} --- posição 6 ($29$) tem pai na posição 3 ($28$) e $29>28$. (b) $33\,32\,28\,31\,29\,26\,25\,30\,27$: \textbf{sim} (todas as 8 relações pai--filho valem).
  \item \emph{Troca com $i<j$, $s_i<s_j$ gera heap?} Falso: $[10,9,1,8,$\allowbreak$7,0,0,6]$, troca das posições 3 e 4 (§\ref{sec:heapprops}).
  \item \emph{Idem com $s_i>s_j$.} Falso: $[10,9]$, troca das posições 1 e 2.
  \item \emph{Número de permutações que são min-heap.} $H(n)=\binom{n-1}{L}H(L)H(n-1-L)$; valores $1,1,2,3,8,$\allowbreak$20,80,210,\dots$ (§\ref{sec:heapprops}).
  \item \emph{Construção linear de $18\,25\,41\,34\,14\,10\,52\,50\,48$.} $[52,50,41,48,$\allowbreak$14,10,18,34,25]$ (traço em §\ref{sec:heap}).
  \item \emph{No máximo $\ceil{n/2^j}$ nós de altura $j$.} Lema \ref{lema:alturas}.
  \item \emph{Construir heap binomial em $O(n)$.} $n$ inserções sucessivas: o custo total é o do contador binário, $\sum_k n/2^k\le2n$ links (§\ref{sec:binomial}). Alternativa: formar pares $B_0+B_0\to B_1$, depois pares de $B_1$, etc., como uma soma binária de $n$ uns.
  \item \emph{Alturas para o pior caso da inserção num heap com $k$ árvores.} Quando o novo $B_0$ provoca ``vai um'' em todas as posições: o heap contém \textbf{$B_0,B_1,\dots,B_{k-1}$} ($n=2^k-1$, todos os bits 1), isto é, alturas $0,1,\dots,k-1$ em arestas ($1,\dots,k$ em níveis, convenção do Szwarcfiter). São feitos $k$ links e o resultado é uma única $B_k$.
  \item[10.] \emph{Remoção do máximo no pior caso.} A remoção custa percorrer as $k$ raízes, gerar os $h$ filhos da raiz removida e fundir $H'$ (as outras $k-1$ árvores) com $H''=\{B_0,\dots,B_{h-1}\}$. O máximo de links ocorre quando $H'$ também contém $B_0,\dots,B_{h-1}$ (cada posição gera ``vai um'') e $h$ é o maior possível: \textbf{$H=\{B_0,B_1,\dots,B_{k-1}\}$ com o máximo na raiz de $B_{k-1}$}, logo $h=k-1$ (em arestas).
  \item[11.] \emph{$k\ge\min\{k_1,k_2\}$?} \textbf{Falso}: $H_1$ com $3=11_2$ nós ($k_1=2$), $H_2$ com $5=101_2$ ($k_2=2$); a fusão tem $8=1000_2$ nós, $k=1<2$.
  \item[12.] \emph{$k\le k_1+k_2$?} \textbf{Verdadeiro}: a união inicial tem $k_1+k_2$ árvores e cada link transforma duas árvores em uma; logo $k=k_1+k_2-\#\text{links}\le k_1+k_2$.
  \item[13.] \emph{Altura mínima e máxima de heap de Fibonacci em função do posto $k$.} Hipótese do exercício (Szwarcfiter §9.5): árvores de Fibonacci de posto $k$, formadas por subordinações de árvores de mesmo posto e cortes controlados por marcas em nós \emph{não raiz}. \textbf{Mínima}: pelo lema $\operatorname{grau}(y_i)\ge i-2$, a menor altura $A(k)$ satisfaz $A(k)=1+A(k-2)$, $A(0)=0$, $A(1)=1$, logo \textbf{$A(k)=\ceil{k/2}$} (atingida pelas árvores de Fibonacci mínimas). \textbf{Máxima}: \textbf{$k$}, atingida pela própria $B_k$ (por indução: subordinar duas árvores de posto $h$ e altura $\le h$ dá posto $h+1$ e altura $\le h+1$; cortes não aumentam a altura nem mudam o posto da raiz). Alturas em arestas.
  (Observação fora do escopo do exercício: numa implementação geral, em que a raiz pode perder filhos livremente, a altura não é limitada por uma função do posto --- o que é sempre $O(\log n)$ é o posto/grau.)
  \item[14--15.] \emph{Algoritmos dos heaps binomiais e de Fibonacci.} §\ref{sec:binomial} e §\ref{sec:fib}.
  \item[16.] \emph{Heap binário sem recursão.}
\end{enumerate}
\begin{pseudo}{Subir e descer iterativos (ex.\ 16; Szwarcfiter ex.\ 6.7--6.8)}
\begin{algorithmic}
\Procedure{subir}{$i$}
  \While{$i>1$ \textbf{e} $T[i].chave>T[\floor{i/2}].chave$}
    \State $T[i]\Leftrightarrow T[\floor{i/2}]$;\ $i:=\floor{i/2}$
  \EndWhile
\EndProcedure
\Procedure{descer}{$i,n$}
  \While{$2i\le n$}
    \State $j:=2i$
    \If{$j<n$ \textbf{e} $T[j+1].chave>T[j].chave$} \State $j:=j+1$ \EndIf
    \If{$T[i].chave<T[j].chave$} \State $T[i]\Leftrightarrow T[j]$;\ $i:=j$
    \Else \State \textbf{pare}
    \EndIf
  \EndWhile
\EndProcedure
\end{algorithmic}
\end{pseudo}
Inserção, remoção, alteração e construção usam exatamente os mesmos esquemas da §\ref{sec:heap} com essas versões.

% ------------------------------------------------------------------
\subsection{Lista -- Union-Find e AVL}
\begin{enumerate}
  \item União por rank sem compressão: altura $\le\floor{\log n}$ e sequência que atinge $\Theta(\log n)$ $\Rightarrow$ Teo.~\ref{teo:ufaltura} e exemplos da §\ref{sec:uf}.
  \item Só compressão: sequência que gera árvore alta e custo $O(m\log n)$ $\Rightarrow$ §\ref{sec:uf}.
  \item Remoção em AVL com rotação em todos os ancestrais $\Rightarrow$ árvore de Fibonacci $T_7$, remover a folha 33 (§\ref{sec:avl}).
  \item Fusão de duas AVL $\Rightarrow$ junção pela espinha direita $O(\log n)$, ou intercalação dos percursos em ordem $O(n+m)$ (§\ref{sec:avl}).
  \item AVL como conjuntos disjuntos: \Find{} $O(\log n)$, \Union{} $O(k\log n)$ $\Rightarrow$ §\ref{sec:avl}.
\end{enumerate}

% ------------------------------------------------------------------
\subsection{Monitoria de revisão: respostas curtas}
\begin{itemize}
  \item \textbf{UF-1.} Objetivo: manter partição em conjuntos disjuntos com consultas ``mesmo conjunto?'' e uniões. \MakeSet{} cria $\{x\}$ ($pai[x]=x$); \Find{} sobe até a raiz; \Union{} liga uma raiz na outra. Cada elemento armazena o \textbf{pai} (e a raiz, o rank). União por rank: liga o de menor rank sob o de maior; altura $\le\log n$.
  \item \textbf{AVL-1.} Balanceada: altura $O(\log n)$ em toda subárvore. AVL: $|h_D-h_E|\le1$ em todo nó. Toda AVL é balanceada (Teo.~\ref{teo:avlaltura}), nem toda balanceada é AVL. Fator de balanço $\bal(v)=h_D(v)-h_E(v)$.
  \item \textbf{AVL-2} (inserção à esquerda de $p$): $\bal=1\to0$, altura não muda, para; $\bal=0\to-1$, altura aumenta, continua; $\bal=-1\to-2$, rotação e para.
  \item \textbf{AVL-3.} $\bal(p)=-2$, $\bal(u)=-1$: simples à direita. $\bal(p)=-2$, $\bal(u)=+1$: dupla à direita (esquerda em $u$, direita em $p$). Diferença: no primeiro o excesso está ``em linha'' (esquerda-esquerda) e uma rotação resolve; no segundo está em ``zigue-zague'' (esquerda-direita) e é preciso primeiro alinhar.
  \item \textbf{Heap-1.} Definição $s_i\le s_{\floor{i/2}}$; seleção $O(1)$ porque o máximo está na posição 1 (Prop.~\ref{prop:ancestral}); custos: $O(1)$, $O(\log n)$, $O(\log n)$, $O(\log n)$, $O(n)$.
  \item \textbf{Binomial (Q6--Q8, Q11).} $B_k$: duas $B_{k-1}$ ligadas; $2^k$ nós; grau da raiz $k$; o heap é a lista de raízes com ordens distintas (bits de $n$). Link de duas $B_2$ de raízes 10 e 15 (min-heap): raiz 10, a de 15 vira sua filha; resultado $B_3$ com 8 nós e grau 3. Mínimo: percorre as $O(\log n)$ raízes (ou ponteiro). União $O(\log n)$: só percorre as listas de raízes (não os $n$ nós), fazendo links como numa soma binária.
  \item \textbf{Fibonacci (Q9, Q10, Q12).} Diferença estrutural: permite várias árvores de mesma ordem (preguiçoso), listas circulares duplas, marcas. $O(1)$\am: inserção, união, aumento de prioridade (\emph{decrease-key} no min-heap), seleção. Remoção do extremo: $O(\log n)$\am. Aumentar prioridade pode violar a ordem com o pai; corta-se $x$ e o leva à lista de raízes; cortes em cascata sobem pelos pais marcados; custo $O(1)$\am\ porque cada corte em cascata é pago pela operação que marcou o nó (potencial $t+2m$).
\end{itemize}

\begin{chave}[Checklist final para a AV1]
\begin{enumerate}
  \item \textbf{AVL}: inserir/remover desenhando a árvore e os balanços a cada passo; rotacionar no desregulado \emph{mais profundo}; saber os 4 casos e o pseudocódigo (papel de $h$ e de $pt$).
  \item \textbf{Union-Find}: tabela de implementações e cenários (listar membros, rollback, offline, Kruskal); provar altura $\le\log n$ com rank.
  \item \textbf{Heap}: provas por caminho/transitividade, índices e indução; construção $O(n)$; contraexemplos pequenos.
  \item \textbf{Escolha}: identificar a operação dominante e o tipo da chave; justificar com complexidade, memória, previsibilidade (pior caso vs.\ amortizado), estabilidade.
  \item \textbf{D\&C}: caso base, divisão, duas chamadas, combinação do ``cruzamento'', recorrência resolvida.
\end{enumerate}
\end{chave}


\end{document}
